% =============================================================================
%  progress3.tex — รายงานความก้าวหน้าโครงการ ULMs ครั้งที่ 3
%
%  build:  cd ~/Projects/SoftwareEn/docs && latexmk progress3.tex
%          ผลลัพธ์ออกที่ build/progress3.pdf (ตาม $out_dir ใน .latexmkrc)
%
%  ต่อจาก progress2.tex (ครั้งที่ 2 · 10 ส.ค. 2569) ซึ่งยังเก็บไว้ไม่แก้
%
%  สิ่งที่ต่างจากสองฉบับแรก
%    1) ลดข้อความเชิงความเห็นลง ทุกข้อความสถานะผูกกับหลักฐานที่ตรวจสอบได้
%       (path ของไฟล์ · จำนวนบรรทัด · commit · หมายเลข CI run) ดูภาคผนวก ค
%    2) เส้น "แผน" ในกราฟ S-Curve ถูกสร้างใหม่จากแผน Sprint คูณน้ำหนักมิติงาน
%       แทนการวาดด้วยมือ และวางจุดตามขอบ Sprint ไม่ใช่ตามสัปดาห์ปฏิทิน
%       (§4.1) โดยยังคงพล็อตเส้นแผนเดิมของรายงานครั้งที่ 1–2 ไว้ให้เทียบด้วย
%    3) เพิ่มกราฟแท่งเทียบรายมิติ ตามสไลด์ Week 06 หน้า 23
%
%  ตัวเลขทุกตัวอ้างสถานะของ origin/main ที่ commit e76f01f ณ 13 ส.ค. 2569
%  โดย "ไม่นับ" โฟลเดอร์ spike-option-b/ และไม่นับโค้ดที่ Prisma สร้างอัตโนมัติ
%  ใน backend/src/generated/
%
%  เขียนรายงานครั้งถัดไป: แก้ "ตัวแปรของรอบรายงาน" ด้านล่าง แล้วไล่แก้
%  §3 (ตารางมิติงาน) · §4 (จุดบนกราฟ) · §5 · §10 · ภาคผนวก ก ข ค
% =============================================================================
% =============================================================================
%  ulms-preamble.tex — preamble ที่ใช้ร่วมกันทุกเอกสารของโครงการ ULMs
%
%  ใช้โดย:  main.tex (ข้อเสนอโครงการ) · progress.tex (รายงานความก้าวหน้า)
%
%  เอกสารที่เรียกใช้ต้องทำสองอย่างหลัง % =============================================================================
%  ulms-preamble.tex — preamble ที่ใช้ร่วมกันทุกเอกสารของโครงการ ULMs
%
%  ใช้โดย:  main.tex (ข้อเสนอโครงการ) · progress.tex (รายงานความก้าวหน้า)
%
%  เอกสารที่เรียกใช้ต้องทำสองอย่างหลัง % =============================================================================
%  ulms-preamble.tex — preamble ที่ใช้ร่วมกันทุกเอกสารของโครงการ ULMs
%
%  ใช้โดย:  main.tex (ข้อเสนอโครงการ) · progress.tex (รายงานความก้าวหน้า)
%
%  เอกสารที่เรียกใช้ต้องทำสองอย่างหลัง \input{ulms-preamble}:
%    1) \hypersetup{pdftitle={...}}  — ชื่อเอกสารใน metadata ของ PDF
%    2) \begin{document}
%  (pdftitle ถูกย้ายออกจากไฟล์นี้ เพราะเป็นค่าเฉพาะของแต่ละเอกสาร)
%
%  แก้ที่นี่ที่เดียว = มีผลกับทุกเอกสาร ห้ามก๊อป preamble ไปวางซ้ำ
% =============================================================================
% =============================================================================
%  main.tex  —  เอกสารข้อเสนอโครงการ (Academic Project Proposal)
%  ระบบบริหารจัดการการยืม–คืนอุปกรณ์มหาวิทยาลัย (ULMs)
%
%  IMPORTANT: This document contains Thai text and MUST be compiled with
%  XeLaTeX (not pdflatex). Thai needs a Unicode font + ICU line-breaking.
%      Build:  latexmk main.tex        (uses .latexmkrc -> xelatex)
%      or:     xelatex main.tex        (run twice for the table of contents)
% =============================================================================

\documentclass[11pt, a4paper]{article}

% ---- Fonts & Thai support ----------------------------------------------------
% fontspec lets XeLaTeX use any system font. Laksaman is a complete serif face
% that covers BOTH Thai and Latin glyphs, so one \setmainfont handles the whole
% (Thai + English) document without font switching.
\usepackage{fontspec}

% Thai has no spaces between words, so TeX cannot find line-break points on its
% own. These two XeTeX primitives switch on ICU's dictionary-based Thai word
% segmentation, which is what makes Thai paragraphs wrap correctly.
\XeTeXlinebreaklocale "th"
\XeTeXlinebreakskip = 0pt plus 0.1pt

\setmainfont{Laksaman}

% ---- Page geometry & paragraph style ----------------------------------------
\usepackage[margin=2.5cm]{geometry}
\usepackage{parskip}          % block paragraphs (no indent, gap between) — common
                              % and more readable for Thai documents
\linespread{1.15}             % a little extra leading; Thai has tall tone marks

% ---- Colours (monochrome / grayscale palette) ------------------------------
\usepackage[table]{xcolor}    % 'table' option enables \rowcolor for table headers
% Monochrome (grayscale) palette. Names kept as-is so every reference still
% works; only the tones changed from the original green/blue accents.
\definecolor{kugreen}{HTML}{1A1A1A}   % near-black — section headings & title
\definecolor{kudark}{HTML}{333333}    % dark gray — subsection/subsubsection titles
\definecolor{headbg}{HTML}{EAEAEA}    % light gray background for table header rows
\definecolor{linkblue}{HTML}{404040}  % medium-dark gray — links / citations

% ---- Section heading styling -------------------------------------------------
% 'nobottomtitles*' stops a heading from being printed near the bottom of a
% page: if less than \bottomtitlespace of room is left, the heading (and its
% following text) moves to the next page. This is titlesec's built-in, robust
% replacement for the earlier \needspace hack — the latter injected vertical
% glue that clashed with longtable page-splitting ("Infinite glue shrinkage").
\usepackage[nobottomtitles*]{titlesec}
\titleformat{\section}
  {\color{kugreen}\Large\bfseries}{\thesection}{0.6em}{}
\titleformat{\subsection}
  {\color{kudark}\large\bfseries}{\thesubsection}{0.5em}{}
\titleformat{\subsubsection}
  {\color{kudark}\normalsize\bfseries}{\thesubsubsection}{0.5em}{}
\titlespacing*{\section}{0pt}{1.4\baselineskip}{0.5\baselineskip}

% Minimum room a heading needs at the bottom of a page; if only ~1–3 lines are
% left (less than this), the heading is bumped to the next page instead of being
% stranded above its table/text.
\renewcommand{\bottomtitlespace}{4\baselineskip}

% ---- Hierarchical indentation by heading level ------------------------------
% Each heading level shifts BOTH the heading and the body text under it further
% to the right: section = 0, subsection (x.y) = 1 step, subsubsection (x.y.z) =
% 2 steps. We do this by having the sectioning commands set \leftskip (a global
% left margin that carries through every following line until the next heading
% resets it). \par first, so the *previous* paragraph is not re-indented.
% \leftskip indents plain paragraphs; \@totalleftmargin makes list environments
% (itemize/enumerate) start from the same indent instead of the page margin.
\makeatletter
\newlength{\lvlindent}
\setlength{\lvlindent}{1.8em}          % one indent step
\let\ulmsSecIndent\section
\renewcommand{\section}{\par\global\leftskip=0pt\global\@totalleftmargin=0pt\relax\ulmsSecIndent}
\let\ulmsSubsecIndent\subsection
\renewcommand{\subsection}{\par\global\leftskip=\lvlindent\global\@totalleftmargin=\lvlindent\relax\ulmsSubsecIndent}
\let\ulmsSubsubsecIndent\subsubsection
\renewcommand{\subsubsection}{\par\global\leftskip=2\lvlindent\global\@totalleftmargin=2\lvlindent\relax\ulmsSubsubsecIndent}
\makeatother

% Indent the headings themselves to match (titlesec ignores \leftskip for the
% heading line, so its left margin is set here via the first \titlespacing arg).
\titlespacing*{\subsection}   {\lvlindent} {1.0\baselineskip}{0.4\baselineskip}
\titlespacing*{\subsubsection}{2\lvlindent}{0.8\baselineskip}{0.3\baselineskip}

% ---- Lists (tighter, custom bullets) ----------------------------------------
\usepackage{enumitem}
\setlist[itemize]{leftmargin=1.6em, itemsep=2pt, topsep=3pt}
\setlist[enumerate]{leftmargin=1.9em, itemsep=2pt, topsep=3pt}

% ---- Tables ------------------------------------------------------------------
\usepackage{array}
\usepackage{booktabs}         % \toprule / \midrule / \bottomrule
\usepackage{tabularx}         % auto-width columns that wrap long Thai text
\usepackage{longtable}        % tables that can break across pages
\usepackage{multirow}         % merged cells for the team-structure table

% Left-aligned wrapping column types:
%   L{width} — fixed-width ragged-right paragraph column
%   Y        — flexible ragged-right X column (for tabularx)
\newcolumntype{L}[1]{>{\raggedright\arraybackslash}p{#1}}
\newcolumntype{Y}{>{\raggedright\arraybackslash}X}
\renewcommand{\arraystretch}{1.3}

% Helpers for a styled header row inside tables.
%   \hrow  — starts a header row and shades it (must be FIRST in the row;
%            colortbl allows only one \rowcolor per row).
%   \thead — bold header-cell text.
\newcommand{\hrow}{\rowcolor{headbg}}
\newcommand{\thead}[1]{\textbf{#1}}

% \begin{keeptable} … \end{keeptable} — keeps a table and the bold label above
% it on the SAME page (used by the data dictionary in §5). A tabularx cannot
% split across pages, so LaTeX is free to break between the label and its
% table, stranding the label alone at the bottom of a page; wrapping the pair
% in a minipage makes them one unbreakable block.
%   * \leftskip is zeroed inside: the minipage box is already placed at the
%     current heading indent, so leaving it set would indent the content twice
%     and push the table past the right margin.
%   * \parskip is restored because minipage's \@parboxrestore zeroes it, which
%     would glue the table straight onto the label line.
% Tables inside must use \linewidth (= the minipage width), not \textwidth.
\newenvironment{keeptable}
  {\par\medskip\noindent
   \begin{minipage}{\dimexpr\textwidth-\leftskip\relax}%
   \leftskip=0pt\relax
   \setlength{\parskip}{0.5\baselineskip plus 2pt}}
  {\end{minipage}\par\medskip}

% ---- Table of contents styling ----------------------------------------------
% By default LaTeX gives dotted leaders to subsections but NOT to sections
% (they print bold with just a right-aligned page number). tocloft lets us turn
% on dots for the section level too, so every TOC line runs "….." to its page.
\usepackage{tocloft}
\renewcommand{\cftsecdotsep}{\cftdotsep}   % dotted leaders for section entries
\renewcommand{\cftsecleader}{\cftdotfill{\cftsecdotsep}}
\renewcommand{\cftsecpagefont}{\normalfont}  % keep page numbers un-bold

% ---- Figures -----------------------------------------------------------------
\usepackage{graphicx}
\graphicspath{{figures/}}

% Landscape (rotated) full-page floats — needed by the ER diagram in §5, which
% is far wider than it is tall and would be unreadable scaled to \textwidth.
\usepackage{rotating}

% ---- Flowcharts / diagrams (TikZ) -------------------------------------------
% Used for the borrow–return workflow diagram (§5). Libraries: geometric shapes
% for decision diamonds, arrows.meta for arrowheads, positioning for relative
% node placement, calc/fit/backgrounds for routing loops and grouping.
\usepackage{tikz}
\usetikzlibrary{shapes.geometric, arrows.meta, positioning, calc, fit, backgrounds}

% ---- Links & bookmarks -------------------------------------------------------
\usepackage{hyperref}
\hypersetup{
  colorlinks=true, linkcolor=black, citecolor=linkblue, urlcolor=linkblue,
  pdflang={th},
}

% Thai labels for auto-generated headings.
\renewcommand{\contentsname}{สารบัญ (Table of Contents)}
\renewcommand{\tablename}{ตาราง}
\renewcommand{\figurename}{รูปที่}

:
%    1) \hypersetup{pdftitle={...}}  — ชื่อเอกสารใน metadata ของ PDF
%    2) \begin{document}
%  (pdftitle ถูกย้ายออกจากไฟล์นี้ เพราะเป็นค่าเฉพาะของแต่ละเอกสาร)
%
%  แก้ที่นี่ที่เดียว = มีผลกับทุกเอกสาร ห้ามก๊อป preamble ไปวางซ้ำ
% =============================================================================
% =============================================================================
%  main.tex  —  เอกสารข้อเสนอโครงการ (Academic Project Proposal)
%  ระบบบริหารจัดการการยืม–คืนอุปกรณ์มหาวิทยาลัย (ULMs)
%
%  IMPORTANT: This document contains Thai text and MUST be compiled with
%  XeLaTeX (not pdflatex). Thai needs a Unicode font + ICU line-breaking.
%      Build:  latexmk main.tex        (uses .latexmkrc -> xelatex)
%      or:     xelatex main.tex        (run twice for the table of contents)
% =============================================================================

\documentclass[11pt, a4paper]{article}

% ---- Fonts & Thai support ----------------------------------------------------
% fontspec lets XeLaTeX use any system font. Laksaman is a complete serif face
% that covers BOTH Thai and Latin glyphs, so one \setmainfont handles the whole
% (Thai + English) document without font switching.
\usepackage{fontspec}

% Thai has no spaces between words, so TeX cannot find line-break points on its
% own. These two XeTeX primitives switch on ICU's dictionary-based Thai word
% segmentation, which is what makes Thai paragraphs wrap correctly.
\XeTeXlinebreaklocale "th"
\XeTeXlinebreakskip = 0pt plus 0.1pt

\setmainfont{Laksaman}

% ---- Page geometry & paragraph style ----------------------------------------
\usepackage[margin=2.5cm]{geometry}
\usepackage{parskip}          % block paragraphs (no indent, gap between) — common
                              % and more readable for Thai documents
\linespread{1.15}             % a little extra leading; Thai has tall tone marks

% ---- Colours (monochrome / grayscale palette) ------------------------------
\usepackage[table]{xcolor}    % 'table' option enables \rowcolor for table headers
% Monochrome (grayscale) palette. Names kept as-is so every reference still
% works; only the tones changed from the original green/blue accents.
\definecolor{kugreen}{HTML}{1A1A1A}   % near-black — section headings & title
\definecolor{kudark}{HTML}{333333}    % dark gray — subsection/subsubsection titles
\definecolor{headbg}{HTML}{EAEAEA}    % light gray background for table header rows
\definecolor{linkblue}{HTML}{404040}  % medium-dark gray — links / citations

% ---- Section heading styling -------------------------------------------------
% 'nobottomtitles*' stops a heading from being printed near the bottom of a
% page: if less than \bottomtitlespace of room is left, the heading (and its
% following text) moves to the next page. This is titlesec's built-in, robust
% replacement for the earlier \needspace hack — the latter injected vertical
% glue that clashed with longtable page-splitting ("Infinite glue shrinkage").
\usepackage[nobottomtitles*]{titlesec}
\titleformat{\section}
  {\color{kugreen}\Large\bfseries}{\thesection}{0.6em}{}
\titleformat{\subsection}
  {\color{kudark}\large\bfseries}{\thesubsection}{0.5em}{}
\titleformat{\subsubsection}
  {\color{kudark}\normalsize\bfseries}{\thesubsubsection}{0.5em}{}
\titlespacing*{\section}{0pt}{1.4\baselineskip}{0.5\baselineskip}

% Minimum room a heading needs at the bottom of a page; if only ~1–3 lines are
% left (less than this), the heading is bumped to the next page instead of being
% stranded above its table/text.
\renewcommand{\bottomtitlespace}{4\baselineskip}

% ---- Hierarchical indentation by heading level ------------------------------
% Each heading level shifts BOTH the heading and the body text under it further
% to the right: section = 0, subsection (x.y) = 1 step, subsubsection (x.y.z) =
% 2 steps. We do this by having the sectioning commands set \leftskip (a global
% left margin that carries through every following line until the next heading
% resets it). \par first, so the *previous* paragraph is not re-indented.
% \leftskip indents plain paragraphs; \@totalleftmargin makes list environments
% (itemize/enumerate) start from the same indent instead of the page margin.
\makeatletter
\newlength{\lvlindent}
\setlength{\lvlindent}{1.8em}          % one indent step
\let\ulmsSecIndent\section
\renewcommand{\section}{\par\global\leftskip=0pt\global\@totalleftmargin=0pt\relax\ulmsSecIndent}
\let\ulmsSubsecIndent\subsection
\renewcommand{\subsection}{\par\global\leftskip=\lvlindent\global\@totalleftmargin=\lvlindent\relax\ulmsSubsecIndent}
\let\ulmsSubsubsecIndent\subsubsection
\renewcommand{\subsubsection}{\par\global\leftskip=2\lvlindent\global\@totalleftmargin=2\lvlindent\relax\ulmsSubsubsecIndent}
\makeatother

% Indent the headings themselves to match (titlesec ignores \leftskip for the
% heading line, so its left margin is set here via the first \titlespacing arg).
\titlespacing*{\subsection}   {\lvlindent} {1.0\baselineskip}{0.4\baselineskip}
\titlespacing*{\subsubsection}{2\lvlindent}{0.8\baselineskip}{0.3\baselineskip}

% ---- Lists (tighter, custom bullets) ----------------------------------------
\usepackage{enumitem}
\setlist[itemize]{leftmargin=1.6em, itemsep=2pt, topsep=3pt}
\setlist[enumerate]{leftmargin=1.9em, itemsep=2pt, topsep=3pt}

% ---- Tables ------------------------------------------------------------------
\usepackage{array}
\usepackage{booktabs}         % \toprule / \midrule / \bottomrule
\usepackage{tabularx}         % auto-width columns that wrap long Thai text
\usepackage{longtable}        % tables that can break across pages
\usepackage{multirow}         % merged cells for the team-structure table

% Left-aligned wrapping column types:
%   L{width} — fixed-width ragged-right paragraph column
%   Y        — flexible ragged-right X column (for tabularx)
\newcolumntype{L}[1]{>{\raggedright\arraybackslash}p{#1}}
\newcolumntype{Y}{>{\raggedright\arraybackslash}X}
\renewcommand{\arraystretch}{1.3}

% Helpers for a styled header row inside tables.
%   \hrow  — starts a header row and shades it (must be FIRST in the row;
%            colortbl allows only one \rowcolor per row).
%   \thead — bold header-cell text.
\newcommand{\hrow}{\rowcolor{headbg}}
\newcommand{\thead}[1]{\textbf{#1}}

% \begin{keeptable} … \end{keeptable} — keeps a table and the bold label above
% it on the SAME page (used by the data dictionary in §5). A tabularx cannot
% split across pages, so LaTeX is free to break between the label and its
% table, stranding the label alone at the bottom of a page; wrapping the pair
% in a minipage makes them one unbreakable block.
%   * \leftskip is zeroed inside: the minipage box is already placed at the
%     current heading indent, so leaving it set would indent the content twice
%     and push the table past the right margin.
%   * \parskip is restored because minipage's \@parboxrestore zeroes it, which
%     would glue the table straight onto the label line.
% Tables inside must use \linewidth (= the minipage width), not \textwidth.
\newenvironment{keeptable}
  {\par\medskip\noindent
   \begin{minipage}{\dimexpr\textwidth-\leftskip\relax}%
   \leftskip=0pt\relax
   \setlength{\parskip}{0.5\baselineskip plus 2pt}}
  {\end{minipage}\par\medskip}

% ---- Table of contents styling ----------------------------------------------
% By default LaTeX gives dotted leaders to subsections but NOT to sections
% (they print bold with just a right-aligned page number). tocloft lets us turn
% on dots for the section level too, so every TOC line runs "….." to its page.
\usepackage{tocloft}
\renewcommand{\cftsecdotsep}{\cftdotsep}   % dotted leaders for section entries
\renewcommand{\cftsecleader}{\cftdotfill{\cftsecdotsep}}
\renewcommand{\cftsecpagefont}{\normalfont}  % keep page numbers un-bold

% ---- Figures -----------------------------------------------------------------
\usepackage{graphicx}
\graphicspath{{figures/}}

% Landscape (rotated) full-page floats — needed by the ER diagram in §5, which
% is far wider than it is tall and would be unreadable scaled to \textwidth.
\usepackage{rotating}

% ---- Flowcharts / diagrams (TikZ) -------------------------------------------
% Used for the borrow–return workflow diagram (§5). Libraries: geometric shapes
% for decision diamonds, arrows.meta for arrowheads, positioning for relative
% node placement, calc/fit/backgrounds for routing loops and grouping.
\usepackage{tikz}
\usetikzlibrary{shapes.geometric, arrows.meta, positioning, calc, fit, backgrounds}

% ---- Links & bookmarks -------------------------------------------------------
\usepackage{hyperref}
\hypersetup{
  colorlinks=true, linkcolor=black, citecolor=linkblue, urlcolor=linkblue,
  pdflang={th},
}

% Thai labels for auto-generated headings.
\renewcommand{\contentsname}{สารบัญ (Table of Contents)}
\renewcommand{\tablename}{ตาราง}
\renewcommand{\figurename}{รูปที่}

:
%    1) \hypersetup{pdftitle={...}}  — ชื่อเอกสารใน metadata ของ PDF
%    2) \begin{document}
%  (pdftitle ถูกย้ายออกจากไฟล์นี้ เพราะเป็นค่าเฉพาะของแต่ละเอกสาร)
%
%  แก้ที่นี่ที่เดียว = มีผลกับทุกเอกสาร ห้ามก๊อป preamble ไปวางซ้ำ
% =============================================================================
% =============================================================================
%  main.tex  —  เอกสารข้อเสนอโครงการ (Academic Project Proposal)
%  ระบบบริหารจัดการการยืม–คืนอุปกรณ์มหาวิทยาลัย (ULMs)
%
%  IMPORTANT: This document contains Thai text and MUST be compiled with
%  XeLaTeX (not pdflatex). Thai needs a Unicode font + ICU line-breaking.
%      Build:  latexmk main.tex        (uses .latexmkrc -> xelatex)
%      or:     xelatex main.tex        (run twice for the table of contents)
% =============================================================================

\documentclass[11pt, a4paper]{article}

% ---- Fonts & Thai support ----------------------------------------------------
% fontspec lets XeLaTeX use any system font. Laksaman is a complete serif face
% that covers BOTH Thai and Latin glyphs, so one \setmainfont handles the whole
% (Thai + English) document without font switching.
\usepackage{fontspec}

% Thai has no spaces between words, so TeX cannot find line-break points on its
% own. These two XeTeX primitives switch on ICU's dictionary-based Thai word
% segmentation, which is what makes Thai paragraphs wrap correctly.
\XeTeXlinebreaklocale "th"
\XeTeXlinebreakskip = 0pt plus 0.1pt

\setmainfont{Laksaman}

% ---- Page geometry & paragraph style ----------------------------------------
\usepackage[margin=2.5cm]{geometry}
\usepackage{parskip}          % block paragraphs (no indent, gap between) — common
                              % and more readable for Thai documents
\linespread{1.15}             % a little extra leading; Thai has tall tone marks

% ---- Colours (monochrome / grayscale palette) ------------------------------
\usepackage[table]{xcolor}    % 'table' option enables \rowcolor for table headers
% Monochrome (grayscale) palette. Names kept as-is so every reference still
% works; only the tones changed from the original green/blue accents.
\definecolor{kugreen}{HTML}{1A1A1A}   % near-black — section headings & title
\definecolor{kudark}{HTML}{333333}    % dark gray — subsection/subsubsection titles
\definecolor{headbg}{HTML}{EAEAEA}    % light gray background for table header rows
\definecolor{linkblue}{HTML}{404040}  % medium-dark gray — links / citations

% ---- Section heading styling -------------------------------------------------
% 'nobottomtitles*' stops a heading from being printed near the bottom of a
% page: if less than \bottomtitlespace of room is left, the heading (and its
% following text) moves to the next page. This is titlesec's built-in, robust
% replacement for the earlier \needspace hack — the latter injected vertical
% glue that clashed with longtable page-splitting ("Infinite glue shrinkage").
\usepackage[nobottomtitles*]{titlesec}
\titleformat{\section}
  {\color{kugreen}\Large\bfseries}{\thesection}{0.6em}{}
\titleformat{\subsection}
  {\color{kudark}\large\bfseries}{\thesubsection}{0.5em}{}
\titleformat{\subsubsection}
  {\color{kudark}\normalsize\bfseries}{\thesubsubsection}{0.5em}{}
\titlespacing*{\section}{0pt}{1.4\baselineskip}{0.5\baselineskip}

% Minimum room a heading needs at the bottom of a page; if only ~1–3 lines are
% left (less than this), the heading is bumped to the next page instead of being
% stranded above its table/text.
\renewcommand{\bottomtitlespace}{4\baselineskip}

% ---- Hierarchical indentation by heading level ------------------------------
% Each heading level shifts BOTH the heading and the body text under it further
% to the right: section = 0, subsection (x.y) = 1 step, subsubsection (x.y.z) =
% 2 steps. We do this by having the sectioning commands set \leftskip (a global
% left margin that carries through every following line until the next heading
% resets it). \par first, so the *previous* paragraph is not re-indented.
% \leftskip indents plain paragraphs; \@totalleftmargin makes list environments
% (itemize/enumerate) start from the same indent instead of the page margin.
\makeatletter
\newlength{\lvlindent}
\setlength{\lvlindent}{1.8em}          % one indent step
\let\ulmsSecIndent\section
\renewcommand{\section}{\par\global\leftskip=0pt\global\@totalleftmargin=0pt\relax\ulmsSecIndent}
\let\ulmsSubsecIndent\subsection
\renewcommand{\subsection}{\par\global\leftskip=\lvlindent\global\@totalleftmargin=\lvlindent\relax\ulmsSubsecIndent}
\let\ulmsSubsubsecIndent\subsubsection
\renewcommand{\subsubsection}{\par\global\leftskip=2\lvlindent\global\@totalleftmargin=2\lvlindent\relax\ulmsSubsubsecIndent}
\makeatother

% Indent the headings themselves to match (titlesec ignores \leftskip for the
% heading line, so its left margin is set here via the first \titlespacing arg).
\titlespacing*{\subsection}   {\lvlindent} {1.0\baselineskip}{0.4\baselineskip}
\titlespacing*{\subsubsection}{2\lvlindent}{0.8\baselineskip}{0.3\baselineskip}

% ---- Lists (tighter, custom bullets) ----------------------------------------
\usepackage{enumitem}
\setlist[itemize]{leftmargin=1.6em, itemsep=2pt, topsep=3pt}
\setlist[enumerate]{leftmargin=1.9em, itemsep=2pt, topsep=3pt}

% ---- Tables ------------------------------------------------------------------
\usepackage{array}
\usepackage{booktabs}         % \toprule / \midrule / \bottomrule
\usepackage{tabularx}         % auto-width columns that wrap long Thai text
\usepackage{longtable}        % tables that can break across pages
\usepackage{multirow}         % merged cells for the team-structure table

% Left-aligned wrapping column types:
%   L{width} — fixed-width ragged-right paragraph column
%   Y        — flexible ragged-right X column (for tabularx)
\newcolumntype{L}[1]{>{\raggedright\arraybackslash}p{#1}}
\newcolumntype{Y}{>{\raggedright\arraybackslash}X}
\renewcommand{\arraystretch}{1.3}

% Helpers for a styled header row inside tables.
%   \hrow  — starts a header row and shades it (must be FIRST in the row;
%            colortbl allows only one \rowcolor per row).
%   \thead — bold header-cell text.
\newcommand{\hrow}{\rowcolor{headbg}}
\newcommand{\thead}[1]{\textbf{#1}}

% \begin{keeptable} … \end{keeptable} — keeps a table and the bold label above
% it on the SAME page (used by the data dictionary in §5). A tabularx cannot
% split across pages, so LaTeX is free to break between the label and its
% table, stranding the label alone at the bottom of a page; wrapping the pair
% in a minipage makes them one unbreakable block.
%   * \leftskip is zeroed inside: the minipage box is already placed at the
%     current heading indent, so leaving it set would indent the content twice
%     and push the table past the right margin.
%   * \parskip is restored because minipage's \@parboxrestore zeroes it, which
%     would glue the table straight onto the label line.
% Tables inside must use \linewidth (= the minipage width), not \textwidth.
\newenvironment{keeptable}
  {\par\medskip\noindent
   \begin{minipage}{\dimexpr\textwidth-\leftskip\relax}%
   \leftskip=0pt\relax
   \setlength{\parskip}{0.5\baselineskip plus 2pt}}
  {\end{minipage}\par\medskip}

% ---- Table of contents styling ----------------------------------------------
% By default LaTeX gives dotted leaders to subsections but NOT to sections
% (they print bold with just a right-aligned page number). tocloft lets us turn
% on dots for the section level too, so every TOC line runs "….." to its page.
\usepackage{tocloft}
\renewcommand{\cftsecdotsep}{\cftdotsep}   % dotted leaders for section entries
\renewcommand{\cftsecleader}{\cftdotfill{\cftsecdotsep}}
\renewcommand{\cftsecpagefont}{\normalfont}  % keep page numbers un-bold

% ---- Figures -----------------------------------------------------------------
\usepackage{graphicx}
\graphicspath{{figures/}}

% Landscape (rotated) full-page floats — needed by the ER diagram in §5, which
% is far wider than it is tall and would be unreadable scaled to \textwidth.
\usepackage{rotating}

% ---- Flowcharts / diagrams (TikZ) -------------------------------------------
% Used for the borrow–return workflow diagram (§5). Libraries: geometric shapes
% for decision diamonds, arrows.meta for arrowheads, positioning for relative
% node placement, calc/fit/backgrounds for routing loops and grouping.
\usepackage{tikz}
\usetikzlibrary{shapes.geometric, arrows.meta, positioning, calc, fit, backgrounds}

% ---- Links & bookmarks -------------------------------------------------------
\usepackage{hyperref}
\hypersetup{
  colorlinks=true, linkcolor=black, citecolor=linkblue, urlcolor=linkblue,
  pdflang={th},
}

% Thai labels for auto-generated headings.
\renewcommand{\contentsname}{สารบัญ (Table of Contents)}
\renewcommand{\tablename}{ตาราง}
\renewcommand{\figurename}{รูปที่}



\hypersetup{pdftitle={รายงานความก้าวหน้าโครงการ ULMs ครั้งที่ 3}}

% ---- การจัดบรรทัด ------------------------------------------------------------
% ภาษาไทยไม่มีการตัดคำด้วยยัติภังค์ และรายงานฉบับนี้มีชื่อไฟล์กับชื่อสัญลักษณ์
% ใน \texttt{} จำนวนมากซึ่ง TeX ตัดกลางคำไม่ได้ ทำให้บางบรรทัดล้นขอบขวา
% \emergencystretch สั่งให้ TeX ยอมยืดช่องไฟเพิ่มเป็นทางเลือกสุดท้าย
% แทนการปล่อยให้ล้น
\setlength{\emergencystretch}{3em}

% \zb — จุดตัดบรรทัดกว้างศูนย์ สำหรับแทรกในชื่อยาว ๆ ที่ \texttt ตัดเองไม่ได้
\newcommand{\zb}{\hspace{0pt}}

% บีบระยะระหว่างบรรทัดในสารบัญ เอกสารนี้มี 32 รายการซึ่งเกินหน้าเดียวไปบรรทัดเดียว
\setlength{\cftbeforesecskip}{1.6pt}
\setlength{\cftbeforesubsecskip}{0.4pt}

% ---- ตัวแปรของรอบรายงาน — แก้ที่นี่ที่เดียวเวลาออกรายงานรอบใหม่ ----------------
\newcommand{\reportno}{3}
\newcommand{\reportdate}{13 สิงหาคม 2569}
\newcommand{\reportperiod}{24 กรกฎาคม – 13 สิงหาคม 2569}
\newcommand{\deltaperiod}{11 – 13 สิงหาคม 2569}
\newcommand{\actualpct}{39\%}      % ผลจริงถ่วงน้ำหนัก ณ วันที่รายงาน
\newcommand{\planrevpct}{42\%}     % แผนที่ปรับปรุงแล้ว (มติรายงานครั้งที่ 2)
\newcommand{\planbasepct}{53\%}    % แผนเดิมที่พล็อตไว้ในรายงานครั้งที่ 1–2
\newcommand{\prevpct}{32\%}        % ผลจริงรายงานครั้งที่ 2
\newcommand{\headsha}{\texttt{e76f01f}}

\begin{document}

% =============================================================================
%  TITLE PAGE
% =============================================================================
\begin{titlepage}
\centering
\vspace*{2.5cm}

{\Large\bfseries\color{kugreen} รายงานความก้าวหน้าโครงการ ครั้งที่ \reportno\par}
\vspace{0.6cm}
{\LARGE\bfseries\color{kugreen} ระบบบริหารจัดการการยืม–คืนอุปกรณ์มหาวิทยาลัย\par}
\vspace{0.35cm}
{\Large\color{kudark} University Lending Management System (ULMs)\par}

\vspace{1.6cm}
\rule{0.72\textwidth}{0.8pt}
\vspace{0.9cm}

\begin{tabular}{@{}r@{\hspace{1.2em}}l@{}}
\textbf{รอบรายงาน}       & \reportperiod \\[0.35em]
\textbf{วันที่รายงาน}     & \reportdate \\[0.35em]
\textbf{สถานะ Sprint}    & Sprint 1 (วันสุดท้าย · วันที่ 14 จาก 14) \\[0.35em]
\textbf{ความก้าวหน้า}    & \actualpct\ \ (ครั้งที่ 1 = 24\% · ครั้งที่ 2 = \prevpct) \\[0.35em]
\textbf{เทียบแผน}        & \planrevpct\ ตามแผนที่ปรับปรุงแล้ว · \planbasepct\ ตามแผนเดิม \\[0.35em]
\textbf{ฐานข้อมูลอ้างอิง} & \texttt{origin/main} ที่ commit \headsha \\
\end{tabular}

\vspace{0.9cm}
\rule{0.72\textwidth}{0.8pt}

\vfill
{\large\color{kudark} คณะทำงาน ``miro ปั่น 14 บาท''\par}
\vspace{0.3cm}
{\color{kudark} ภาควิชาวิศวกรรมคอมพิวเตอร์ คณะวิศวกรรมศาสตร์\par}
{\color{kudark} มหาวิทยาลัยเกษตรศาสตร์\par}
\vspace{1.2cm}
\end{titlepage}

\tableofcontents
\clearpage

% =============================================================================
\section{บทสรุปผู้บริหาร (Executive Summary)}

ณ วันที่ \reportdate\ ซึ่งเป็นวันสุดท้ายของ Sprint 1 โครงการใช้เวลาไปแล้ว 2.9 สัปดาห์จากแผนทั้งหมด 10 สัปดาห์ คิดเป็น 29\% ของระยะเวลา ความก้าวหน้าของงานถ่วงน้ำหนักอยู่ที่ \actualpct\ และค่าใช้จ่ายอยู่ที่ 0 บาทจากวงเงินประเมิน 5{,}000 บาท

\subsection*{ตัวเลขสามชุดที่ต้องอ่านคู่กัน}

\begin{tabularx}{\dimexpr\textwidth-\leftskip\relax}{@{} L{5.6cm} L{2.0cm} L{2.0cm} Y @{}}
\toprule
\hrow \thead{เทียบกับ} & \thead{แผน} & \thead{ผลจริง} & \thead{ผลต่าง} \\
\midrule
แผนเดิม (พล็อตไว้ในรายงานครั้งที่ 1–2) & \planbasepct & \actualpct & $-14$ จุด \\
แผนที่ปรับปรุงแล้ว (มติรายงานครั้งที่ 2) & \planrevpct & \actualpct & $-3$ จุด \\
ผลจริงของรายงานครั้งที่แล้ว & --- & \prevpct & $+7$ จุด \\
\bottomrule
\end{tabularx}

\vspace{0.8em}
\noindent ผลต่าง $-14$ จุดเทียบแผนเดิมเป็นค่าเดียวกับที่รายงานไว้ในครั้งที่ 2 ส่วนผลต่าง $-3$ จุดเทียบแผนที่ปรับปรุงแล้วนั้นเป็นการเทียบกับแผนที่คณะทำงานเสนอไว้ในรายงานครั้งที่ 2 หัวข้อ 10 วิธีสร้างเส้นแผนทั้งสองเส้นอธิบายไว้ในหัวข้อ \ref{sec:curvemethod} และค่าที่ใช้พล็อตอยู่ในภาคผนวก ข

\subsection*{ผลงานที่เพิ่มขึ้นในช่วง \deltaperiod}

นับเฉพาะที่เข้าสาย \texttt{main} แล้ว และไม่นับโค้ดที่ Prisma สร้างอัตโนมัติ รวม 53 ไฟล์ เพิ่ม 5{,}365 บรรทัด ลบ 292 บรรทัด แยกเป็น

\begin{itemize}
  \item \textbf{ฝั่ง Back-End} 9 ไฟล์ 471 บรรทัด เป็น tRPC Procedure แรกของโครงการ (\texttt{auth.me}) พร้อมชั้นควบคุมสิทธิ์ 4 บทบาท ชั้น Context และไฟล์ทดสอบชุดแรก
  \item \textbf{ฝั่ง Front-End} 40 ไฟล์ เพิ่ม 4{,}465 บรรทัด เป็นหน้าจอผู้ยืม 3 หน้า หน้าโปรไฟล์ 1 หน้า และส่วนประกอบร่วม
  \item \textbf{ฝั่งฐานข้อมูล} ปรับสคีมาใหม่เป็น 27 ตารางกับ 9 Enum และสร้าง Migration ใหม่ \texttt{20260813063940\zb\_init}
  \item \textbf{ฝั่งเครื่องมือ} รวมระบบตรวจสอบอัตโนมัติ (CI) เข้าสาย \texttt{main} แล้ว
\end{itemize}

\subsection*{สามข้อเท็จจริงที่กำหนดสถานะของ Sprint 2}

\begin{enumerate}
  \item \textbf{มี tRPC Procedure ในระบบแล้ว 1 รายการ จากที่สัญญาระบุไว้ 60 รายการ} คือ \texttt{auth.me} ในไฟล์ \texttt{backend/src/auth/auth.router.ts}
  \item \textbf{Procedure ดังกล่าวยังเรียกสำเร็จไม่ได้} เพราะฟังก์ชัน \texttt{verifySessionToken()} ใน \texttt{backend/src/trpc/context.ts} บรรทัดที่ 82–84 คืนค่า \texttt{null} เสมอ (มีคอมเมนต์ \texttt{TODO} กำกับ) ทำให้ \texttt{ctx.user} เป็น \texttt{null} ทุกครั้ง และมิดเดิลแวร์จะคืนค่า \texttt{UNAUTHORIZED} และยังไม่มี Procedure สำหรับเข้าสู่ระบบ
  \item \textbf{ยังไม่มีไฟล์สัญญากลางในรีโป} ปัจจุบันมีสัญญาสองชุดที่เขียนด้วยมือแยกกันคนละฝั่ง และทั้งสองฝั่งใช้ไลบรารี \texttt{zod} คนละรุ่นหลัก (Back-End \texttt{\^{}4.4.3} · Front-End \texttt{\^{}3.23.8}) สคริปต์ \texttt{scripts/check-versions.mjs} รายงานปัญหา 3 ข้อและคืนค่า exit code 1
\end{enumerate}

\subsection*{ข้อเสนอที่ขอมติจากคณะกรรมการ}

\begin{enumerate}
  \item ล็อกสัญญา tRPC และคอมมิตไฟล์สัญญาเข้ารีโป พร้อมยกรุ่น \texttt{zod} ฝั่ง Front-End จาก 3 เป็น 4 ในคำขอรวมโค้ดเดียวกัน กำหนดภายใน 15 สิงหาคม 2569
  \item อนุมัติการตัดขอบเขต 3 รายการตามที่เสนอไว้ในรายงานครั้งที่ 2 ซึ่งยังไม่มีมติ (หัวข้อ \ref{sec:scopechange})
  \item เพิ่มกำลังคนฝั่ง Back-End ปัจจุบันมีผู้คอมมิตโค้ดฝั่ง Back-End 2 คน ต่อสัดส่วนน้ำหนักงาน 30\% ซึ่งสูงที่สุดในโครงการ
  \item คำแนะนำเรื่องการคงไว้ซึ่ง MongoDB ซึ่งเป็นข้อค้างมาตั้งแต่รายงานครั้งที่ 1
\end{enumerate}

% =============================================================================
\clearpage
\section{สถานะโครงการโดยรวม (เงิน · งาน · เวลา)}

\begin{tabularx}{\dimexpr\textwidth-\leftskip\relax}{@{} L{1.5cm} L{3.3cm} L{2.6cm} Y @{}}
\toprule
\hrow \thead{ด้าน} & \thead{แผน} & \thead{ผลจริง} & \thead{ข้อเท็จจริงประกอบ} \\
\midrule
เวลา & ผ่านไป 2.9 จาก 10 สัปดาห์ (29\%) & ผ่านไป 2.9 จาก 10 สัปดาห์ (29\%) & วันสิ้นสุดโครงการยังเป็น 1 ตุลาคม 2569 ไม่มีการเลื่อน \\
งาน & \planrevpct\ (แผนปรับปรุง) \newline \planbasepct\ (แผนเดิม) & \actualpct & คำนวณถ่วงน้ำหนักตามตารางหัวข้อ \ref{sec:dimension} ที่มาของแต่ละตัวเลขอยู่ในภาคผนวก ก \\
เงิน & 5{,}000 บาท (วงเงินประเมินทั้งโครงการ) & 0 บาท (0\%) & ยังไม่มีรายจ่าย รายการที่ประเมินไว้คือค่าเช่า Cloud/Hosting และฐานข้อมูลกลาง ซึ่งยังไม่ได้จัดหา \\
\bottomrule
\end{tabularx}

\vspace{0.8em}
\noindent วงเงิน 5{,}000 บาทที่ประเมินไว้ในข้อเสนอโครงการผูกกับการนำระบบขึ้นใช้งานจริงและฐานข้อมูลกลางสำหรับทดสอบร่วมกัน ปัจจุบันสมาชิกแต่ละคนใช้ฐานข้อมูล PostgreSQL บนเครื่องตนเอง ค่าใช้จ่าย 0 บาทจึงสอดคล้องกับข้อเท็จจริงที่ว่ายังไม่มีการจัดหาทรัพยากรทั้งสองรายการ

% =============================================================================
\section{ความก้าวหน้าแยกตามมิติของงาน}
\label{sec:dimension}

การแบ่งมิติงานเป็นไปตามสไลด์ Week 06 หน้า 23 คือแยกเป็นส่วนออกแบบ (Front-End Design และ Back-End Design) ส่วนพัฒนา (Front-End Dev และ Back-End Dev) ส่วนทดสอบ และส่วนเอกสาร ค่าน้ำหนักของแต่ละมิติไม่เปลี่ยนจากรายงานครั้งที่ 1 และ 2 เพื่อให้เทียบสามรอบกันได้โดยตรง ที่มาของทุกตัวเลขในคอลัมน์ ``จริง'' แจกแจงไว้ในภาคผนวก ก

\begin{tabularx}{\dimexpr\textwidth-\leftskip\relax}{@{} L{3.0cm} L{1.0cm} L{1.1cm} L{0.9cm} L{1.1cm} L{1.7cm} Y @{}}
\toprule
\hrow \thead{มิติงาน} & \thead{น้ำหนัก} & \thead{แผน} & \thead{จริง} & \thead{ผลต่าง} & \thead{ครั้งที่ 1 / 2} & \thead{ข้อเท็จจริงประกอบ} \\
\midrule
Back-End Design \newline (API + Database Model) & 10\% & 88\% & 64\% & $\mathbf{-24}$ & 45 / 55 & แบบจำลองข้อมูล Migrate ได้จริง สัญญา API ยังไม่มีไฟล์กลางในรีโป \\
Front-End Design \newline (UI/UX + Business Process) & 10\% & 75\% & 76\% & $+1$ & 60 / 72 & ออกแบบหน้าจอที่มีเนื้อหาจริงแล้ว 11 จาก 28 หน้า \\
Back-End Dev & 30\% & 25\% & 20\% & $-5$ & 12 / 12 & มี Procedure 1 จาก 60 รายการ ยังไม่มีสคริปต์ข้อมูลตั้งต้น \\
Front-End Dev & 25\% & 42\% & 46\% & $\mathbf{+4}$ & 20 / 38 & เพิ่มหน้าผู้ยืม 3 หน้าและหน้าโปรไฟล์ ทุกหน้ายังอ่านข้อมูลจำลอง \\
Test (Front-End + Back-End) & 15\% & 10\% & 10\% & $0$ & 0 / 0 & มีไฟล์ทดสอบชุดแรก 1 ไฟล์ 2 กรณี และสคริปต์ตรวจเกณฑ์ยอมรับ 1 ไฟล์ \\
Documentation & 10\% & 65\% & 64\% & $-1$ & 45 / 62 & เพิ่มแบบฟอร์มลงมติสัญญา ยังไม่มี SRS และ SDS ฉบับทางการ \\
\midrule
\hrow \thead{รวมถ่วงน้ำหนัก} & \thead{100\%} & \thead{\planrevpct} & \thead{\actualpct} & \thead{$\mathbf{-3}$} & \thead{24 / 32} & \\
\bottomrule
\end{tabularx}

\vspace{0.8em}
\noindent คอลัมน์ ``แผน'' ในตารางนี้คือค่าเป้าหมายของแผนที่ปรับปรุงแล้ว ณ วันสิ้นสุด Sprint 1 ตามที่เสนอไว้ในรายงานครั้งที่ 2 หัวข้อ 10.1 ไม่ใช่ค่าของแผนเดิม การเทียบกับแผนเดิมอยู่ในกราฟหัวข้อ \ref{sec:scurve}

ผลต่างรวม $-3$ จุดมาจากสองมิติที่ต่ำกว่าแผนและสองมิติที่สูงกว่าแผน แจกแจงเป็นจุดถ่วงน้ำหนักได้ดังนี้ Back-End Design $-2.4$ จุด · Back-End Dev $-1.5$ จุด · Documentation $-0.1$ จุด · Front-End Dev $+1.0$ จุด · Front-End Design $+0.1$ จุด รวมเป็น $-2.9$ จุด ปัดเป็น $-3$

\vspace{0.5em}
\noindent เมื่อเทียบสามรอบรายงานติดกัน มิติที่เปลี่ยนแปลงมากที่สุดคือ Front-End Dev (20 $\rightarrow$ 38 $\rightarrow$ 46) และ Back-End Dev ซึ่งขยับเป็นครั้งแรกในรอบนี้ (12 $\rightarrow$ 12 $\rightarrow$ 20) ส่วนมิติ Test ขยับจาก 0 เป็นค่าบวกเป็นครั้งแรกเช่นกัน

% =============================================================================
\clearpage
\section{กราฟความก้าวหน้าเทียบแผน (S-Curve)}
\label{sec:scurve}

\subsection{วิธีสร้างเส้นแผนและเหตุผลที่ปรับ}
\label{sec:curvemethod}

สไลด์ Week 06 หน้า 13 ถึง 21 กำหนดหลักการของกราฟไว้สามข้อ คือต้องพล็อต Planned Progress เทียบ Actual Progress บนแกนเดียวกัน · เส้นแผนต้องมีรูปทรงตัว S คือเริ่มช้า เร่งกลางโครงการ แล้วชะลอตอนท้าย · และต้องอ่านออกได้ทันทีว่างานเร็วหรือช้ากว่าแผน รายงานฉบับนี้ยึดทั้งสามข้อไว้ครบ และปรับเฉพาะ\textbf{วิธีได้มาซึ่งค่าบนเส้นแผน}

\vspace{0.5em}
\noindent\textbf{เส้นแผนเดิม} ที่ใช้ในรายงานครั้งที่ 1 และ 2 คือชุดพิกัด 7 จุดที่กำหนดค่าไว้โดยตรง ได้แก่ $(0,0)$ $(1,25)$ $(3,55)$ $(5,75)$ $(7,90)$ $(9,97)$ $(10,100)$ จุดเหล่านี้วางตามสัปดาห์ปฏิทินและไม่ได้คำนวณย้อนกลับไปหาแผน Sprint หรือค่าน้ำหนักมิติงาน จึงตรวจสอบไม่ได้ว่าค่า 55\% ณ สัปดาห์ที่ 3 มาจากงานรายการใดบ้าง

\vspace{0.5em}
\noindent\textbf{เส้นแผนที่ปรับปรุงแล้ว} ในรายงานฉบับนี้สร้างขึ้นด้วยขั้นตอนสามขั้นที่ตรวจสอบย้อนกลับได้ทุกจุด

\begin{enumerate}
  \item \textbf{กำหนดค่าเป้าหมายรายมิติ ณ วันสิ้นสุดของแต่ละ Sprint} จากรายการงานที่แผน Sprint นั้นระบุไว้ เช่น เมื่อสิ้นสุด Sprint 2 แผนกำหนดให้วงจรยืม–คืนเส้นทางปกติทำงานได้ Back-End Dev จึงตั้งเป้าไว้ที่ 60\%
  \item \textbf{คูณค่าเป้าหมายแต่ละมิติด้วยน้ำหนักของมิตินั้น แล้วรวมกัน} ได้เป็นค่าความก้าวหน้าสะสมที่ควรอยู่ ณ ขอบ Sprint นั้น ตารางค่าทั้งหมดอยู่ในภาคผนวก ข
  \item \textbf{วางจุดบนกราฟที่ขอบ Sprint ไม่ใช่ที่สัปดาห์ปฏิทิน} แล้วลากเส้นเชื่อม
\end{enumerate}

\vspace{0.5em}
\noindent เหตุผลที่ปรับให้เป็นแบบนี้มีสามข้อ ซึ่งล้วนมาจากรูปแบบการทำงานจริงของกลุ่ม

\begin{tabularx}{\dimexpr\textwidth-\leftskip\relax}{@{} L{4.4cm} Y @{}}
\toprule
\hrow \thead{ลักษณะการทำงานของกลุ่ม} & \thead{การปรับที่ทำกับเส้นแผน} \\
\midrule
เกณฑ์ ``เสร็จ'' ของกลุ่มถูกประเมินที่ขอบ Sprint เท่านั้น ไม่ได้ประเมินรายสัปดาห์ & ย้ายจุดบนเส้นแผนไปวางที่ขอบ Sprint ทั้ง 6 จุด แทนการวางที่สัปดาห์ที่ 1, 3, 5, 7, 9 ตามเดิม แกนนอนจึงเป็นขอบ Sprint และมีเส้นประกำกับทุกขอบ \\
งานเข้าสาย \texttt{main} เป็นก้อนตอนรวมโค้ด ไม่ได้ทยอยเข้าทุกวัน ในรอบนี้มีการรวมโค้ด 5 ครั้งภายใน 3 วัน & เส้นแผนเป็นเส้นตรงต่อกันระหว่างขอบ Sprint ไม่ทำให้โค้งเรียบ เพราะความชันระหว่างขอบสองขอบคือปริมาณงานที่ Sprint นั้นรับไว้ ซึ่งเป็นค่าที่มีความหมายจริงกับกลุ่ม \\
Sprint 3 (28 ส.ค. – 10 ก.ย.) ทับช่วงสอบกลางภาค กำลังคนที่ใช้ได้จริงลดลง & ลดความชันของช่วง Sprint 3 ลงเหลือ 6.5 จุดต่อสัปดาห์ ต่ำที่สุดในบรรดา Sprint ที่ยังเพิ่มความสามารถใหม่ แล้วไปเร่งคืนใน Sprint 4 เส้นแผนจึงมีรอยหยักตรงช่วงสอบ ซึ่งเส้นแผนเดิมไม่มี \\
\bottomrule
\end{tabularx}

\vspace{0.8em}
\noindent ผลของการปรับคือความชันของเส้นแผนไล่จาก 14 · 15 · 11.5 · \textbf{6.5} · 8 · 6 จุดต่อสัปดาห์ ตั้งแต่ Sprint 0 ถึง Sprint 5 ซึ่งเป็นรูปทรงตัว S ตามหลักการในสไลด์ คือไต่ขึ้นสูงสุดช่วงกลาง แล้วชะลอลงตอนท้าย โดยมีรอยหยักช่วงสอบกลางภาคเพิ่มเข้ามา

\vspace{0.5em}
\noindent \textbf{เส้นแผนเดิมยังคงถูกพล็อตไว้ในกราฟด้วย} เพื่อให้ตรวจสอบได้ว่าการเปลี่ยนวิธีคำนวณทำให้ค่าเปลี่ยนไปเท่าใด และเพื่อไม่ให้การปรับวิธีคำนวณกลบผลต่างที่รายงานไว้เดิม

\subsection{กราฟ S-Curve}

\begin{figure}[!ht]
\centering
\begin{tikzpicture}[x=1.20cm, y=0.058cm]
  % ---- กริดแนวนอน ----
  \foreach \y in {0,20,40,60,80,100}
    \draw[gray!22] (0,\y) -- (10.2,\y);

  % ---- เส้นประกำกับขอบ Sprint ----
  \foreach \x in {0.86, 2.86, 4.86, 6.86, 8.86}
    \draw[gray!40, dashed, thin] (\x,0) -- (\x,108);

  % ---- ป้ายชื่อ Sprint ด้านบน ----
  \foreach \x/\lbl in {0.43/{S0}, 1.86/{S1}, 3.86/{S2}, 5.86/{S3}, 7.86/{S4}, 9.36/{S5}}
    \node[font=\scriptsize\bfseries, gray!50!black] at (\x,114) {\lbl};

  % ---- แกน ----
  \draw[->, thick] (0,0) -- (10.6,0);
  \draw[->, thick] (0,0) -- (0,122) node[above, font=\scriptsize] {\% ความก้าวหน้าสะสม};
  \foreach \y in {0,20,40,60,80,100}
    \draw (0,\y) -- (-0.10,\y) node[left, font=\scriptsize] {\y};

  % ---- ขีดและวันที่ตามขอบ Sprint ----
  \foreach \x/\lbl in {0/{24 ก.ค.}, 0.86/{30 ก.ค.}, 2.86/{13 ส.ค.},
                       4.86/{27 ส.ค.}, 6.86/{10 ก.ย.}, 8.86/{24 ก.ย.}, 9.86/{1 ต.ค.}}
    {\draw (\x,0) -- (\x,-3);
     \node[rotate=45, anchor=north east, font=\scriptsize] at (\x,-4) {\lbl};}

  % ---- เส้นแผนเดิม (รายงานครั้งที่ 1–2) ----
  \draw[thick, gray!55, dashed]
    plot coordinates {(0,0) (1,25) (3,55) (5,75) (7,90) (9,97) (10,100)};

  % ---- เส้นแผนที่ปรับปรุงแล้ว — จุดวางที่ขอบ Sprint ----
  \draw[thick, kudark, densely dashed]
    plot coordinates {(0,1) (0.86,12) (2.86,42) (4.86,65) (6.86,78) (8.86,94) (9.86,100)};
  \foreach \p in {(0,1), (0.86,12), (2.86,42), (4.86,65), (6.86,78), (8.86,94), (9.86,100)}
    \fill[kudark] \p circle (1.5pt);

  % ---- เส้นผลจริง — จุดคือวันที่ออกรายงานแต่ละครั้ง ----
  \draw[very thick, kugreen]
    plot coordinates {(0,0) (2.0,24) (2.43,32) (2.86,39)};
  \foreach \p in {(0,0), (2.0,24), (2.43,32), (2.86,39)}
    \fill[kugreen] \p circle (2.2pt);

  % ---- ช่องว่างเทียบแผนสองเส้น ณ วันที่รายงาน ----
  % ลูกศรวางเยื้องขวาของจุดรายงาน เพื่อไม่ให้ทับตัวจุดเอง
  \draw[<->, red!70!black, thick] (3.06,39) -- (3.06,53);
  \node[right, font=\scriptsize, red!70!black, align=left] at (3.18,48)
    {ต่ำกว่าแผนเดิม 14 จุด};
  \node[right, font=\scriptsize, kudark, align=left] at (3.18,36)
    {ต่ำกว่าแผนปรับปรุง 3 จุด};

  % ---- รายการจุดบนเส้นผลจริง วางในมุมบนซ้ายซึ่งเป็นพื้นที่ว่างของกราฟ ----
  % fill=white เพื่อบังเส้นกริดที่พาดผ่านข้อความ
  \node[anchor=north west, font=\scriptsize, align=left, gray!25!black,
        fill=white, inner sep=4pt] at (0.22,105)
    {จุดบนเส้นผลจริง\\[2pt]
     {\scriptsize ครั้งที่ 1 · 7 ส.ค. · 24\%}\\
     {\scriptsize ครั้งที่ 2 · 10 ส.ค. · 32\%}\\
     {\scriptsize \textbf{ครั้งที่ 3 · 13 ส.ค. · 39\%}}};

  % ---- คำอธิบายสัญลักษณ์ (บังกริดด้วยพื้นขาวก่อนวาด) ----
  \fill[white] (5.35,2) rectangle (10.2,34);
  \draw[thick, gray!55, dashed] (5.5,28) -- (6.2,28);
  \node[right, font=\scriptsize] at (6.3,28) {แผนเดิม (รายงานครั้งที่ 1–2)};
  \draw[thick, kudark, densely dashed] (5.5,18) -- (6.2,18);
  \node[right, font=\scriptsize] at (6.3,18) {แผนที่ปรับปรุงแล้ว (ฉบับนี้)};
  \draw[very thick, kugreen] (5.5,8) -- (6.2,8);
  \node[right, font=\scriptsize] at (6.3,8) {ผลจริง (Actual)};
\end{tikzpicture}
\caption{กราฟ S-Curve เปรียบเทียบความก้าวหน้าตามแผนกับผลการดำเนินงานจริง
  แกนนอนวางตามขอบ Sprint ทั้ง 6 ช่วง (เส้นประแนวตั้ง)
  เส้นแผนที่ปรับปรุงแล้วมีจุดที่ขอบ Sprint ทุกจุด และมีความชันต่ำที่สุดในช่วง Sprint 3 ซึ่งทับช่วงสอบกลางภาค
  ค่าที่ใช้พล็อตทั้งหมดอยู่ในภาคผนวก ข}
\label{fig:scurve}
\end{figure}

\noindent อ่านจากรูปที่ \ref{fig:scurve} ได้สามข้อ ข้อแรก เส้นผลจริงมีความชัน 16 จุดต่อสัปดาห์ในช่วง 6 วันหลัง (24\% ณ 7 ส.ค. เป็น 39\% ณ 13 ส.ค.) เทียบกับความชันของเส้นแผนที่ปรับปรุงแล้วในช่วง Sprint 1 ซึ่งอยู่ที่ 15 จุดต่อสัปดาห์ ข้อสอง ระยะห่างระหว่างเส้นผลจริงกับเส้นแผนที่ปรับปรุงแล้วแคบลงจาก 11 จุด ณ 7 ส.ค. เหลือ 3 จุด ณ 13 ส.ค. ข้อสาม ระยะห่างระหว่างเส้นผลจริงกับเส้นแผนเดิมยังคงเป็น 14 จุด เท่ากับที่รายงานไว้ในครั้งที่ 2

\clearpage
\subsection{กราฟเทียบรายมิติ ณ วันที่รายงาน}

สไลด์ Week 06 หน้า 23 กำหนดให้แบ่งส่วนงานออกเป็นมิติเพื่อให้เห็นผลงานรายด้าน กราฟแท่งนี้แสดงค่าของทั้ง 6 มิติที่ใช้คำนวณจุด 39\% บนเส้นผลจริง เทียบกับค่าเป้าหมายของแผนที่ปรับปรุงแล้ว ณ วันเดียวกัน

\begin{figure}[!ht]
\centering
\begin{tikzpicture}[x=0.062cm, y=1cm]
  % ---- กริดแนวตั้งและสเกลแกนนอน ----
  \foreach \x in {0,20,40,60,80,100} {
    \draw[gray!25] (\x,-0.30) -- (\x,6.15);
    \node[below, font=\scriptsize] at (\x,-0.30) {\x};
  }
  \node[below, font=\scriptsize] at (50,-0.68) {\% ความก้าวหน้าของมิตินั้น};
  \draw[thick] (0,-0.30) -- (0,6.15);

  % ---- แท่งรายมิติ: y / ชื่อ / น้ำหนัก / แผน / จริง ----
  \foreach \y/\name/\wt/\plan/\act in {%
      5.6/{Back-End Design}/10/88/64,
      4.6/{Front-End Design}/10/75/76,
      3.6/{Back-End Dev}/30/25/20,
      2.6/{Front-End Dev}/25/42/46,
      1.6/{Test}/15/10/10,
      0.6/{Documentation}/10/65/64} {
    \node[left, xshift=-0.16cm, font=\scriptsize, align=right] at (0,\y)
      {\name\\[-3pt]{\tiny น้ำหนัก \wt\%}};
    \fill[gray!45]  (0,{\y+0.05}) rectangle (\plan,{\y+0.30});
    \fill[kugreen]  (0,{\y-0.30}) rectangle (\act,{\y-0.05});
    \node[right, font=\tiny, gray!45!black] at (\plan,{\y+0.175}) {\plan};
    \node[right, font=\tiny, kugreen]       at (\act,{\y-0.175}) {\act};
  }

  % ---- คำอธิบายสัญลักษณ์ — วางสองบรรทัด เพื่อไม่ให้ข้อความสองป้ายชนกัน ----
  \fill[gray!45] (0,-1.36) rectangle (7,-1.54);
  \node[right, font=\scriptsize] at (9,-1.45) {แผนที่ปรับปรุงแล้ว};
  \fill[kugreen] (0,-1.76) rectangle (7,-1.94);
  \node[right, font=\scriptsize] at (9,-1.85) {ผลจริง};
\end{tikzpicture}
\vspace{0.6em}
\caption{ความก้าวหน้ารายมิติ ณ 13 สิงหาคม 2569 เทียบกับค่าเป้าหมายของแผนที่ปรับปรุงแล้ว
  ตัวเลขในวงเล็บใต้ชื่อมิติคือน้ำหนักที่ใช้ถ่วงเมื่อรวมเป็นค่าเดียว
  มิติที่มีน้ำหนักสูงสุดคือ Back-End Dev (30\%) ซึ่งอยู่ที่ 20\%}
\label{fig:dims}
\end{figure}

\noindent รูปที่ \ref{fig:dims} แสดงว่ามิติที่ห่างจากแผนมากที่สุดคือ Back-End Design ($-24$ จุด) ซึ่งมาจากส่วนสัญญา API ที่ยังไม่มีไฟล์กลางในรีโป ส่วนมิติที่มีน้ำหนักสูงสุดคือ Back-End Dev อยู่ที่ 20\% เทียบแผน 25\% และมิติที่สูงกว่าแผนคือ Front-End Dev และ Front-End Design

% =============================================================================
\clearpage
\section{แผนและผลการดำเนินงานในรอบที่ผ่านมา}

\subsection{แผนที่กำหนดไว้สำหรับช่วง \deltaperiod\ และผลที่เกิดขึ้น}

รายการต่อไปนี้คือแผนที่คณะทำงานเสนอไว้ในรายงานครั้งที่ 2 หัวข้อ 10.1 สำหรับสามวันสุดท้ายของ Sprint 1 พร้อมผลที่ตรวจสอบได้จากรีโป

\begin{tabularx}{\dimexpr\textwidth-\leftskip\relax}{@{} L{1.3cm} L{4.4cm} L{1.9cm} Y @{}}
\toprule
\hrow \thead{กำหนด} & \thead{งานตามแผน} & \thead{ผล} & \thead{หลักฐาน} \\
\midrule
11 ส.ค. & สคริปต์ข้อมูลตั้งต้น 2 หน่วยงาน พร้อมผู้ใช้ครบทุกบทบาท & \textbf{ไม่ได้ทำ} & ไม่พบไฟล์ seed ใน \texttt{backend/prisma/} และไม่มีคีย์ \texttt{prisma.seed} ใน \texttt{package.json} \\
12 ส.ค. & เพิ่มตารางผูกชิ้นอุปกรณ์ สถานะอนุมัติบางส่วน และคีย์ไม่ซ้ำบนอีเมล & \textbf{บางส่วน} & สคีมาถูกปรับใหม่เมื่อ 13 ส.ค. เป็น 27 ตาราง 9 Enum · \texttt{ItemIndiv.ResourceKey} เป็น \texttt{@unique} แล้ว จึงผูกชิ้นผ่าน \texttt{Reservations} $\rightarrow$ \texttt{ResourceInfo} ได้ · \texttt{ApproveStatus} ยังมีเพียง 4 ค่า ไม่มีค่าอนุมัติบางส่วน · \texttt{AccountInfo.Email} และ \texttt{UserID} ยังไม่มี \texttt{@unique} \\
13 ส.ค. & ประชุมลงมติสัญญา tRPC แล้วคอมมิตไฟล์สัญญาเข้ารีโป & \textbf{ไม่ได้ทำ} & ไม่มีไฟล์สัญญากลางในรีโป มีเพียงสัญญาเขียนมือสองชุดแยกฝั่ง (รายละเอียดในหัวข้อ \ref{sec:contract}) \\
13 ส.ค. & รวม CI เข้าสายหลัก & \textbf{เสร็จ} & PR \#19 รวมเข้า \texttt{main} เมื่อ 12 ส.ค. · run ล่าสุด 2 ครั้งบน \texttt{main} ผ่านทั้งคู่ (13 ส.ค.) \\
13 ส.ค. & เกณฑ์การยอมรับของวงจรหลัก & \textbf{บางส่วน} & \texttt{frontend/src/lib/\_\_checks\_\_/ac-checks.ts} 116 บรรทัด ครอบคลุมเกณฑ์ยอมรับข้อ 2.4 ถึง 2.7 ยังไม่มีเอกสารแผนการทดสอบ \\
\bottomrule
\end{tabularx}

\subsection{งานที่ส่งมอบในช่วงเดียวกันแต่ไม่ได้อยู่ในแผนสามวันนี้}

\begin{tabularx}{\dimexpr\textwidth-\leftskip\relax}{@{} L{1.1cm} Y L{3.0cm} @{}}
\toprule
\hrow \thead{ส่วนงาน} & \thead{ผลงาน} & \thead{ที่มาในแผน} \\
\midrule
BE & tRPC Procedure \texttt{auth.me} พร้อมชั้น Context ชั้นควบคุมสิทธิ์ 4 บทบาท Zod Schema 2 ไฟล์ Mapper 1 ไฟล์ และไฟล์ทดสอบ 1 ไฟล์ รวม 471 บรรทัด & Sprint 2 ตามแผนปรับปรุง \newline (ดึงขึ้นมาทำก่อน) \\
FE & หน้ารายการอุปกรณ์ 597 บรรทัด · หน้ารายละเอียดอุปกรณ์ 359 บรรทัด · หน้ารายการห้อง 556 บรรทัด & Sprint 2 ตามแผนปรับปรุง \newline (ดึงขึ้นมาทำก่อน) \\
FE & หน้าโปรไฟล์ผู้ใช้ 242 บรรทัด · เมนูค้นหาแบบคีย์ลัด 169 บรรทัด · แถบเมนูล่างสำหรับมือถือ 97 บรรทัด · เมนูการแจ้งเตือน 145 บรรทัด & ไม่มีในแผนฉบับใด \\
FE & ชั้นเชื่อมต่อ tRPC ฝั่งไคลเอนต์ (\texttt{lib/trpc.ts}) และสัญญาเขียนมือฝั่ง Front-End 60 Procedure (\texttt{server/trpc-contract.ts}) & ไม่มีในแผนฉบับใด \\
FE & ไลบรารีกฎธุรกิจ 165 บรรทัด · ข้อความแสดงข้อผิดพลาด 62 บรรทัด · การตรวจไฟล์อัปโหลด 61 บรรทัด · คำแปลไทย–อังกฤษเพิ่มภาษาละ 158 บรรทัด & Sprint 3 ตามแผนปรับปรุง \newline (ดึงขึ้นมาทำก่อน) \\
\bottomrule
\end{tabularx}

\vspace{0.8em}
\noindent สรุปเชิงปริมาณของช่วง \deltaperiod\ คือมีการรวมโค้ดเข้าสาย \texttt{main} 5 ครั้ง (PR \#16 ถึง \#20) เปลี่ยนแปลง 53 ไฟล์ เพิ่ม 5{,}365 บรรทัด ลบ 292 บรรทัด โดยไม่นับโค้ดที่ Prisma สร้างอัตโนมัติ

\subsection{สถานะเกณฑ์ ``เสร็จ'' ของ Sprint 1}

เกณฑ์ที่ปรับใหม่ในรายงานครั้งที่ 2 มีสามข้อ ผลเมื่อสิ้นสุด Sprint 1 คือ

\begin{tabularx}{\dimexpr\textwidth-\leftskip\relax}{@{} L{5.6cm} L{1.8cm} Y @{}}
\toprule
\hrow \thead{เกณฑ์} & \thead{ผล} & \thead{หมายเหตุ} \\
\midrule
สัญญาถูกล็อกและคอมมิตแล้ว & \textbf{ไม่ผ่าน} & ยังไม่มีไฟล์สัญญากลาง \\
ฐานข้อมูลมีข้อมูลตั้งต้น & \textbf{ไม่ผ่าน} & ยังไม่มีสคริปต์ seed \\
CI ทำงานบนสายหลัก & \textbf{ผ่าน} & run 2 ครั้งล่าสุดบน \texttt{main} ผ่านทั้งคู่ \\
\bottomrule
\end{tabularx}

\vspace{0.8em}
\noindent Sprint 1 จึงผ่านเกณฑ์ 1 ข้อจาก 3 ข้อ งานสองข้อที่ไม่ผ่านยกยอดไปเป็นงานแรกของ Sprint 2

% =============================================================================
\clearpage
\section{ผลงานที่ส่งมอบแล้ว}

หัวข้อนี้รายงานเฉพาะผลงานที่ตรวจสอบได้ในรีโปที่ commit \headsha\ ไม่รวมกิจกรรมระหว่างดำเนินงาน ไม่นับโฟลเดอร์ \texttt{spike-option-b/} ซึ่งเป็นงานทดลองสถาปัตยกรรม และไม่นับ 41 ไฟล์ใน \texttt{backend/src/generated/prisma/} ซึ่ง Prisma สร้างขึ้นอัตโนมัติ

\subsection{ฝั่ง Back-End}

\begin{itemize}
  \item \textbf{โครงสร้างพื้นฐาน} NestJS 11 ประกอบกับ Prisma 7, \texttt{nestjs-trpc} 2.13, \texttt{cookie-parser} และ \texttt{ConfigModule} ใน \texttt{app.module.ts} · \texttt{main.ts} ตั้งค่า CORS แบบระบุ Origin เจาะจงพร้อม \texttt{credentials: true} และเปิดใช้ \texttt{cookieParser()}
  \item \textbf{ฐานข้อมูล} \texttt{schema.prisma} นิยาม 27 ตารางกับ 9 Enum และ Migration \texttt{20260813063940\zb\_init} สร้างตารางทั้งหมดได้จริง (538 บรรทัด SQL)
  \item \textbf{ชั้น Context ของ tRPC} (\texttt{src/trpc/context.ts} 84 บรรทัด) อ่านคุกกี้ \texttt{ulms\_session} แล้วประกอบเป็น \texttt{ctx.user} ที่มี \texttt{accountKey}, \texttt{role}, \texttt{facultyKey} และ \texttt{creditScore}
  \item \textbf{ชั้นควบคุมสิทธิ์} (\texttt{src/trpc/auth.middleware.ts} 48 บรรทัด) มี \texttt{AuthMiddleware} สำหรับบังคับให้เข้าสู่ระบบก่อน และมิดเดิลแวร์ตามบทบาทอีก 3 ตัวคือ \texttt{StaffMiddleware}, \texttt{Supervisor\zb Middleware} และ \texttt{AdminMiddleware} ซึ่งสร้างจากฟังก์ชันโรงงานตัวเดียวกัน
  \item \textbf{tRPC Procedure รายการแรก} \texttt{auth.me} (\texttt{src/auth/auth.router.ts} 17 บรรทัด) คืนโปรไฟล์ผู้ใช้พร้อมเพดานการยืมที่คำนวณจากระดับเครดิต โดยบริการเบื้องหลัง (\texttt{auth.service.ts} 63 บรรทัด) เลือกเฉพาะคอลัมน์ที่ต้องใช้และไม่ดึง \texttt{HashedPassword} ออกมา
  \item \textbf{Zod Schema และ Mapper} \texttt{common/schemas/user.schema.ts} 35 บรรทัด · \texttt{common/schemas/status.schema.ts} 55 บรรทัด · \texttt{common/mappers/user.mapper.ts} 45 บรรทัด
  \item \textbf{ไฟล์ทดสอบชุดแรกของโครงการ} \texttt{src/auth/auth.service.spec.ts} 99 บรรทัด 2 กรณีทดสอบ สร้างข้อมูลตั้งต้นขั้นต่ำเข้าฐานข้อมูลจริงแล้วตรวจว่าผลลัพธ์ผ่านการตรวจด้วย Zod Schema และมีค่าตรงตามที่คาด
\end{itemize}

\noindent\textbf{ข้อจำกัดที่ตรวจพบ}

\begin{itemize}
  \item ฟังก์ชัน \texttt{verifySessionToken()} ใน \texttt{context.ts} คืนค่า \texttt{null} เสมอ (มีคอมเมนต์ \texttt{TODO: replace with real session/JWT verification}) ผลคือ \texttt{ctx.user} เป็น \texttt{null} ทุกครั้ง และการเรียก \texttt{auth.me} จะได้ \texttt{UNAUTHORIZED} เสมอ
  \item ยังไม่มี Procedure สำหรับเข้าสู่ระบบและออกจากระบบ จึงยังไม่มีทางสร้างคุกกี้ \texttt{ulms\_session} ได้
  \item ยังไม่มีสคริปต์สร้างข้อมูลตั้งต้น ฐานข้อมูลที่ Migrate แล้วจึงยังว่างเปล่า
  \item \texttt{AccountInfo.Email} และ \texttt{AccountInfo.UserID} ยังไม่มีข้อกำหนด \texttt{@unique} (รายงานไว้ตั้งแต่ครั้งที่ 2)
  \item ไฟล์ทดสอบ \texttt{auth.service.spec.ts} ต้องต่อฐานข้อมูลจริงจึงจะรันได้ CI ในปัจจุบันจึงยังไม่ได้รันไฟล์นี้
\end{itemize}

\subsection{ฝั่ง Front-End}

\begin{itemize}
  \item \textbf{โครงสร้างพื้นฐาน} Vite 5 + React 18 + TypeScript 5.6 + React Router 6 · เส้นทางหน้าจอ 28 เส้นทาง · ชั้นคุ้มกันสิทธิ์ตามบทบาท 4 กลุ่มซ้อนบนเส้นทางเหล่านั้น · โหลดหน้าจอแบบแยกก้อน (lazy) ทุกหน้า
  \item \textbf{ระบบออกแบบ} Tailwind CSS 3.4 พร้อมโทเคนสี · ส่วนประกอบพื้นฐาน 20 ตัว รวมตารางข้อมูล กราฟ กล่องโต้ตอบ แผงเลื่อน และป้ายระดับเครดิต · รองรับไทย–อังกฤษผ่าน i18next และโหมดสว่าง–มืด
  \item \textbf{หน้าจอที่มีเนื้อหาจริง 11 หน้า จาก 28 หน้า} แยกเป็น
    \begin{itemize}
      \item ผู้ดูแลระบบ 6 หน้า — ผู้ใช้ (612 บรรทัด) · บันทึกการตรวจสอบ (353) · แผงสรุป (286) · สถานะระบบ (274) · ตั้งค่าระบบ (252) · รายงาน (229)
      \item ผู้ยืม 3 หน้า (\textbf{ใหม่ในรอบนี้}) — รายการอุปกรณ์ (597) · รายการห้อง (556) · รายละเอียดอุปกรณ์ (359)
      \item ทุกบทบาท 2 หน้า — โปรไฟล์ผู้ใช้ (242 บรรทัด · \textbf{ใหม่ในรอบนี้}) · เข้าสู่ระบบ (124)
    \end{itemize}
  \item \textbf{ชั้นเชื่อมต่อ tRPC} \texttt{lib/trpc.ts} 37 บรรทัด ตั้งค่า \texttt{httpBatchLink} ไปที่ \texttt{/trpc} พร้อม \texttt{credentials: "include"} และ \texttt{TRPCProvider} ถูกประกอบเข้ากับ \texttt{app.tsx} แล้ว
  \item \textbf{ไลบรารีกฎธุรกิจ} \texttt{lib/business-rules.ts} 165 บรรทัด · \texttt{lib/validation.ts} 43 บรรทัด · \texttt{lib/upload-validation.ts} 61 บรรทัด · \texttt{lib/error-messages.ts} 110 บรรทัด
  \item \textbf{สคริปต์ตรวจเกณฑ์การยอมรับ} \texttt{lib/\_\_checks\_\_/ac-checks.ts} 116 บรรทัด ตรวจกฎเรื่องช่วงเวลายืม เวลาตัดคืน 17:00 น. จำนวนช่องจอง T3 พร้อมกัน และการกรองอุปกรณ์ที่ยืมได้
\end{itemize}

\noindent\textbf{ข้อจำกัดที่ตรวจพบ}

\begin{itemize}
  \item ยังไม่มีไฟล์ใดในโฟลเดอร์ \texttt{src/} เรียกใช้ \texttt{useTRPC()} หรือ \texttt{useTRPCClient()} ชั้นเชื่อมต่อ tRPC จึงถูกประกอบไว้แล้วแต่ยังไม่มีผู้ใช้งาน
  \item Hook ที่ดึงข้อมูลมี 4 ตัว โดย 3 ตัว (\texttt{use-equipment-types}, \texttt{use-my-loans}, \texttt{use-rooms}) คืนค่าจากข้อมูลจำลองในโค้ด และอีก 1 ตัว (\texttt{use-me}) เรียกผ่าน \texttt{apiClient} ไปยังเส้นทางแบบ REST \texttt{/auth/me} ซึ่งฝั่ง Back-End ไม่ได้เปิดให้บริการ
  \item หน้าจอทั้ง 11 หน้าที่มีเนื้อหาจริงยังอ่านข้อมูลจากไฟล์จำลอง 3 ไฟล์ คือ \texttt{features/admin/mock-data.ts} (370 บรรทัด) \texttt{features/borrower/mock-data.ts} (347 บรรทัด) และ \texttt{features/notifications/mock-notifications.ts} (77 บรรทัด)
  \item ยังไม่มีตัวรันชุดทดสอบฝั่ง Front-End ติดตั้งอยู่ (ไม่พบ \texttt{vitest} หรือ \texttt{jest} ใน \texttt{package.json} และไม่มีสคริปต์ \texttt{test}) สคริปต์ \texttt{ac-checks.ts} จึงต้องรันด้วยมือผ่าน esbuild
  \item หน้าจอที่ยังเป็นหน้าเปล่า 17 หน้า ครอบคลุมหน้าของเจ้าหน้าที่ทั้ง 7 หน้า หน้าของผู้อนุมัติ 2 หน้า หน้ารายงาน 2 หน้า หน้าของผู้ยืมที่เหลือ 5 หน้า และหน้าตั้งค่าของผู้ดูแลระบบ 1 หน้า
\end{itemize}

\subsection{ฝั่งเครื่องมือและกระบวนการ}

\begin{itemize}
  \item \textbf{ระบบตรวจสอบอัตโนมัติ (CI)} \texttt{.github/workflows/ci.yml} 150 บรรทัด กำหนดงานตรวจ 3 งาน คือตรวจเวอร์ชันไลบรารีร่วม · ติดตั้ง ตรวจชนิดข้อมูล และ Build ฝั่ง Back-End · ติดตั้ง ตรวจชนิดข้อมูล และ Build ฝั่ง Front-End ทำงานทั้งบนคำขอรวมโค้ดและการ Push เข้า \texttt{main}
  \item \textbf{สคริปต์ตรวจเวอร์ชันร่วม} \texttt{scripts/check-versions.mjs} 141 บรรทัด ตรวจ 3 ข้อคือทั้งสองฝั่งประกาศเวอร์ชันเดียวกัน · เวอร์ชันถูกตรึงแบบเจาะจง · และติดตั้งอยู่ชุดเดียวในแต่ละฝั่ง
  \item \textbf{สถานะการทำงานของ CI} รวมเข้า \texttt{main} เมื่อ 12 สิงหาคม 2569 ผ่าน PR \#19 · run 2 ครั้งแรกวันที่ 12 สิงหาคมล้มเหลว · run 2 ครั้งล่าสุดวันที่ 13 สิงหาคมผ่านทั้งคู่ ใช้เวลา 34 และ 40 วินาที
\end{itemize}

\noindent\textbf{ข้อจำกัดที่ตรวจพบ} งานตรวจเวอร์ชันถูกตั้งค่า \texttt{continue-on-error: true} จึงรายงานผลแต่ไม่ขวางการรวมโค้ด สาเหตุที่ต้องตั้งเช่นนี้คือฝั่ง Back-End ใช้ \texttt{zod} รุ่น \texttt{\^{}4.4.3} ส่วนฝั่ง Front-End ใช้รุ่น \texttt{\^{}3.23.8} เมื่อรันสคริปต์ปัจจุบันจะได้ผลว่ามีปัญหา 3 ข้อและคืนค่า exit code 1

\subsection{ฝั่งเอกสารและการออกแบบ}

\begin{itemize}
  \item \textbf{ข้อเสนอโครงการ 42 หน้า} และ\textbf{การออกแบบฐานข้อมูลรอบที่ 4} พร้อมพจนานุกรมข้อมูล
  \item \textbf{คู่มือสัญญา tRPC 31 หน้า} · \textbf{วาระประชุมเรื่องสัญญา 33 หน้า 19 วาระ} · \textbf{โค้ดตัวอย่างผลลัพธ์ของสัญญา 28 ไฟล์}
  \item \textbf{รายงานผลงานทดลองสถาปัตยกรรม 31 หน้า} (\texttt{docs/report/spike-gap.pdf}) ระบุจุดที่ต้องแก้ 19 จุด เรียงตามความรุนแรง P0 ถึง P3
  \item \textbf{แบบฟอร์มลงมติสัญญา} (\texttt{docs/ULMs-tRPC-decision-form.xlsx} · \textbf{ใหม่ในรอบนี้}) สำหรับบันทึกมติของแต่ละวาระในที่ประชุม
  \item \textbf{สารบัญไฟล์ของโครงการ} (\texttt{docs/INDEX.md}) ระบุว่าเอกสารฉบับล่าสุดของแต่ละเรื่องอยู่ที่ไหน
\end{itemize}

\noindent\textbf{ข้อจำกัดที่ตรวจพบ} ยังไม่มีเอกสาร SRS ฉบับทางการ ยังไม่มี SDS นอกเหนือจากบทที่ว่าด้วยฐานข้อมูล และยังไม่มีคู่มือผู้ใช้กับคู่มือติดตั้ง

% =============================================================================
\clearpage
\section{สถานะของสัญญา API}
\label{sec:contract}

หัวข้อนี้แยกออกมาเพราะเป็นรายการเดียวที่ปรากฏเป็นสาเหตุในทั้งสามรายงาน และเป็นเงื่อนไขของงานส่วนใหญ่ใน Sprint 2

\begin{tabularx}{\dimexpr\textwidth-\leftskip\relax}{@{} L{3.6cm} L{2.2cm} Y @{}}
\toprule
\hrow \thead{สิ่งที่มีอยู่} & \thead{ขนาด} & \thead{ข้อเท็จจริง} \\
\midrule
สัญญาฝั่ง Front-End & 188 บรรทัด \newline 60 Procedure & \texttt{frontend/src/server/trpc-contract.ts} เขียนด้วยมือ ครอบคลุม 11 กลุ่มคือ auth, item, loan, reservation, approval, appeal, credit, inspection, notification, report, admin · ในไฟล์มีคอมเมนต์ระบุว่าเป็นของชั่วคราวและให้แทนที่ด้วยชนิดข้อมูลจริงจากฝั่ง Back-End เมื่อมีแล้ว \\
สัญญาฝั่ง Back-End & 90 บรรทัด \newline 1 Procedure & Zod Schema 2 ไฟล์ใน \texttt{backend/src/common/schemas/} และ \texttt{auth.me} 1 รายการ ประกาศ Output Schema ไว้ที่ตัว Procedure \\
เอกสารประกอบการตัดสินใจ & 64 หน้า \newline 28 ไฟล์ตัวอย่าง & คู่มือ 31 หน้า วาระประชุม 33 หน้า 19 วาระ และโค้ดตัวอย่าง 28 ไฟล์ พร้อมแบบฟอร์มลงมติ \\
\bottomrule
\end{tabularx}

\vspace{0.8em}
\noindent\textbf{ช่องว่างที่เหลืออยู่ 3 ข้อ}

\begin{enumerate}
  \item \textbf{ไม่มีแหล่งความจริงเดียว} สัญญาสองชุดข้างต้นถูกเขียนแยกกันและไม่มีกลไกใดตรวจว่าตรงกัน หากฝั่งใดฝั่งหนึ่งเปลี่ยน อีกฝั่งจะไม่รู้จนกว่าจะเรียกใช้งานจริงแล้วพัง
  \item \textbf{รุ่นของ \texttt{zod} ไม่ตรงกัน} Back-End ใช้ \texttt{\^{}4.4.3} (ติดตั้งจริง 4.4.3) Front-End ใช้ \texttt{\^{}3.23.8} (ติดตั้งจริง 3.25.76) ชนิดข้อมูลของ Zod รุ่น 3 กับ 4 เข้ากันไม่ได้ในระดับชนิด การนำ \texttt{AppRouter} จากฝั่ง Back-End มาใช้ฝั่ง Front-End จึงจะตรวจชนิดข้อมูลไม่ผ่าน
  \item \textbf{เวอร์ชันยังไม่ถูกตรึงแบบเจาะจง} ทั้งสองฝั่งประกาศเป็นช่วง (\texttt{\^{}}) สคริปต์ตรวจจึงรายงานเป็นปัญหาเพิ่มอีก 2 ข้อ เพราะการติดตั้งใหม่คนละเวลาอาจได้คนละรุ่นย่อย
\end{enumerate}

\vspace{0.5em}
\noindent ทั้งสามข้อแก้ได้ในคำขอรวมโค้ดเดียว คือยกรุ่น \texttt{zod} ฝั่ง Front-End เป็น 4 พร้อมยกรุ่น \texttt{@hookform/resolvers} จาก 3 เป็น 5 ตรึงทั้งสองฝั่งเป็นเวอร์ชันเจาะจง แล้วส่งออกชนิด \texttt{AppRouter} จากฝั่ง Back-End มาแทนไฟล์สัญญาชั่วคราว เมื่อทำแล้วสามารถถอด \texttt{continue-on-error} ออกจากงานตรวจเวอร์ชันใน CI ได้

% =============================================================================
\clearpage
\section{การเปลี่ยนแปลงขอบเขตโครงการ}
\label{sec:scopechange}

\subsection{รายการที่เสนอตัดในรายงานครั้งที่ 2 และยังไม่มีมติ}

\begin{tabularx}{\dimexpr\textwidth-\leftskip\relax}{@{} L{3.6cm} Y L{2.6cm} @{}}
\toprule
\hrow \thead{รายการ} & \thead{เหตุผลที่เสนอตัด} & \thead{สถานะ} \\
\midrule
การจองสถานที่ (Facility T3) & ระบุไว้ตั้งแต่ข้อเสนอโครงการว่าเป็นงานที่ตัดได้ และไม่มีงานอื่นขึ้นกับมัน & \textbf{รอมติ} \\
การส่งออกรายงาน (Export) & หน้าจอรายงานของผู้ดูแลระบบทำเสร็จแล้ว ใช้แสดงในการนำเสนอแทนได้ & \textbf{รอมติ} \\
การจองล่วงหน้าขั้นสูง & ได้แก่ การผูกชิ้นอัตโนมัติเมื่อใกล้วันรับ การจัดคิวเมื่อจองไม่สำเร็จ และเพดานการต่ออายุที่เทียบกับวันจอง & \textbf{รอมติ} \\
\bottomrule
\end{tabularx}

\vspace{0.8em}
\noindent ข้อสังเกตที่ต้องบันทึก หน้ารายการห้อง (\texttt{room-list-page.tsx} 556 บรรทัด) ถูกพัฒนาขึ้นในรอบนี้ และเป็นหน้าจอฝั่งผู้ใช้ของความสามารถ ``การจองสถานที่'' ซึ่งอยู่ในรายการที่เสนอตัด หากที่ประชุมมีมติตัด หน้าจอนี้จะไม่มีบริการฝั่ง Back-End รองรับ

\subsection{การเปลี่ยนแปลงอื่นที่ต้องบันทึก}

\begin{tabularx}{\dimexpr\textwidth-\leftskip\relax}{@{} L{3.8cm} Y L{2.2cm} @{}}
\toprule
\hrow \thead{เรื่อง} & \thead{การเปลี่ยนแปลง} & \thead{สถานะ} \\
\midrule
โครงสร้างสคีมา & สคีมาถูกปรับใหม่เมื่อ 13 ส.ค. จาก 33 ตารางเหลือ 27 ตาราง และเพิ่ม 9 Enum เพื่อจำกัดค่าที่เป็นไปได้ของสถานะต่าง ๆ Migration เดิมถูกแทนที่ด้วย \texttt{20260813063940\zb\_init} & \textbf{ดำเนินการแล้ว} \\
ตารางผูกชิ้นอุปกรณ์ & \texttt{ItemIndiv.ResourceKey} เป็น \texttt{@unique} จึงผูกชิ้นอุปกรณ์เข้ากับการจองผ่าน \texttt{Reservations} $\rightarrow$ \texttt{ResourceInfo} ได้ ข้อค้างจากรายงานครั้งที่ 1 และ 2 ในเรื่องนี้ถือว่าปิดแล้ว & \textbf{ปิดแล้ว} \\
สถานะ ``อนุมัติบางส่วน'' & Enum \texttt{ApproveStatus} มี 4 ค่า คือ Pending, Approved, Rejected, Canceled ยังไม่มีค่าสำหรับอนุมัติบางส่วน & \textbf{ยังค้าง} \\
คีย์ไม่ซ้ำบนอีเมลและรหัสผู้ใช้ & \texttt{AccountInfo.Email} และ \texttt{UserID} ยังไม่มี \texttt{@unique} ฐานข้อมูลจึงยอมให้สร้างบัญชีที่มีอีเมลซ้ำกันได้ & \textbf{ยังค้าง} \\
หน้ารายการอุปกรณ์และรายละเอียด & แผนปรับปรุงกำหนดไว้ใน Sprint 2 แต่ทำเสร็จแล้วในช่วง Sprint 1 & \textbf{ดึงขึ้นมาแล้ว} \\
ฐานข้อมูล MongoDB & ยังไม่มีข้อสรุป เป็นข้อค้างจากรายงานครั้งที่ 1 & \textbf{รอคำแนะนำ} \\
\bottomrule
\end{tabularx}

% =============================================================================
\clearpage
\section{ปัญหาอุปสรรคและแนวทางแก้ไข}

\begin{tabularx}{\dimexpr\textwidth-\leftskip\relax}{@{} L{3.0cm} Y L{2.3cm} @{}}
\toprule
\hrow \thead{ปัญหา} & \thead{ข้อเท็จจริง ผลกระทบ และแนวทางแก้ไข} & \thead{ผู้รับผิดชอบ / กำหนด} \\
\midrule
สัญญา API ยังไม่ถูกล็อก \newline \textbf{(ค้างจากครั้งที่ 1 และ 2)} &
\textbf{ข้อเท็จจริง} มีสัญญาเขียนมือ 2 ชุดแยกฝั่ง 60 Procedure เทียบ 1 Procedure และรุ่น \texttt{zod} ไม่ตรงกัน \newline
\textbf{ผลกระทบ} งานฝั่ง Back-End ทุกรายการใน Sprint 2 ขึ้นกับข้อนี้ \newline
\textbf{แนวทาง} ประชุมลงมติวาระที่จำเป็นต่อวงจรหลัก ยกรุ่น \texttt{zod} ฝั่ง Front-End และส่งออก \texttt{AppRouter} จากฝั่ง Back-End ในคำขอรวมโค้ดเดียวกัน &
API + SA \newline ภายใน 15 ส.ค. \\

ระบบยืนยันตัวตนยังทำงานไม่ครบวง \newline \textbf{(ใหม่)} &
\textbf{ข้อเท็จจริง} มีชั้นควบคุมสิทธิ์ครบ 4 บทบาทและมี Procedure \texttt{auth.me} แล้ว แต่ \texttt{verifySessionToken()} คืนค่า \texttt{null} เสมอ และยังไม่มี Procedure เข้าสู่ระบบ \newline
\textbf{ผลกระทบ} ไม่มี Procedure ใดเรียกสำเร็จได้ในขณะนี้ และหน้าจอยังต่อของจริงไม่ได้ \newline
\textbf{แนวทาง} เขียน \texttt{auth.login} ที่ออกคุกกี้ \texttt{ulms\_session} และเติมการตรวจโทเคนใน \texttt{verifySessionToken()} ทั้งสองรายการไม่ขึ้นกับมติสัญญา &
BE-dev \newline ภายใน 17 ส.ค. \\

ไม่มีข้อมูลตั้งต้นในฐานข้อมูล \newline \textbf{(ค้างจากครั้งที่ 1 และ 2)} &
\textbf{ข้อเท็จจริง} ไม่พบไฟล์ seed และไม่มีคีย์ \texttt{prisma.seed} ใน \texttt{package.json} \newline
\textbf{ผลกระทบ} ไฟล์ทดสอบที่มีอยู่ต้องสร้างข้อมูลเองทุกครั้ง และ CI ยังรันชุดทดสอบไม่ได้ \newline
\textbf{แนวทาง} เขียนสคริปต์ seed 2 หน่วยงานพร้อมผู้ใช้ครบทุกบทบาท ไม่ขึ้นกับมติสัญญา &
BE-dev \newline ภายใน 16 ส.ค. \\

หน้าจอทั้งหมดยังอ่านข้อมูลจำลอง \newline \textbf{(ค้างจากครั้งที่ 1 และ 2)} &
\textbf{ข้อเท็จจริง} หน้าจอ 11 หน้าที่มีเนื้อหาจริงอ่านจากไฟล์จำลอง 3 ไฟล์ รวม 794 บรรทัด และยังไม่มีไฟล์ใดเรียก \texttt{useTRPC()} \newline
\textbf{ผลกระทบ} จำนวนหน้าที่ต้องแก้เมื่อต่อข้อมูลจริงเพิ่มขึ้นทุกรอบ จาก 7 หน้าเป็น 11 หน้า \newline
\textbf{แนวทาง} เมื่อ \texttt{auth.login} และ \texttt{auth.me} ใช้ได้จริง ให้เปลี่ยนหน้าโปรไฟล์เป็นหน้าแรกที่ต่อของจริง เพื่อพิสูจน์เส้นทางข้อมูลครบสามชั้นก่อนแก้หน้าที่เหลือ &
FE-lead + BE-dev \newline ภายใน 20 ส.ค. \\

งานทดสอบครอบคลุมน้อย \newline \textbf{(ปรับสถานะ)} &
\textbf{ข้อเท็จจริง} มีไฟล์ทดสอบฝั่ง Back-End 1 ไฟล์ 2 กรณี ซึ่งต้องต่อฐานข้อมูลจริง และฝั่ง Front-End ยังไม่มีตัวรันชุดทดสอบติดตั้งอยู่ \newline
\textbf{ผลกระทบ} มิติ Test มีน้ำหนัก 15\% และอยู่ที่ 10\% \newline
\textbf{แนวทาง} ติดตั้ง \texttt{vitest} ฝั่ง Front-End แล้วย้ายสคริปต์ \texttt{ac-checks.ts} เข้าเป็นชุดทดสอบจริง และเพิ่มงานตรวจชุดทดสอบเข้า CI หลังจากมีสคริปต์ seed แล้ว &
QA-lead \newline ภายใน 20 ส.ค. \\

งานตรวจเวอร์ชันใน CI ยังไม่บังคับใช้ \newline \textbf{(ใหม่)} &
\textbf{ข้อเท็จจริง} งานตรวจถูกตั้ง \texttt{continue-on-error: true} และสคริปต์รายงานปัญหา 3 ข้อเมื่อรันในขณะนี้ \newline
\textbf{ผลกระทบ} ความไม่ตรงกันของรุ่นไลบรารีจะไม่ขวางการรวมโค้ด \newline
\textbf{แนวทาง} ถอด \texttt{continue-on-error} ออกในคำขอรวมโค้ดเดียวกับที่ยกรุ่น \texttt{zod} &
FE-lead \newline ภายใน 15 ส.ค. \\
\bottomrule
\end{tabularx}

% =============================================================================
\clearpage
\section{การประเมินความเสี่ยง}

\begin{tabularx}{\dimexpr\textwidth-\leftskip\relax}{@{} L{3.6cm} L{1.3cm} L{1.5cm} Y @{}}
\toprule
\hrow \thead{ความเสี่ยง} & \thead{โอกาส} & \thead{ผลกระทบ} & \thead{ฐานของการประเมินและแนวทางรับมือ} \\
\midrule
Sprint 2 แบกงานสอง Sprint รวมกันแล้วตามไม่ทัน & \textbf{สูง} & \textbf{สูง} & Sprint 2 รับเป้าหมายเดิมของ Sprint 1 ทั้งชุด บวกงานตกค้างจาก Sprint 0 อีก 2 รายการ · กำหนดจุดตัดสินใจกลาง Sprint วันที่ 20 ส.ค. หากยังสร้างคำขอผ่านระบบจริงไม่ได้ ให้คงไว้เฉพาะเส้นทางเดียวที่สาธิตได้ \\
งานกระจุกที่ Sprint 3 ซึ่งทับช่วงสอบกลางภาค & \textbf{สูง} & \textbf{สูง} & Sprint 3 (28 ส.ค. – 10 ก.ย.) รับเนื้อหาเดิมของ Sprint 2 มาทั้งชุดหลังการเลื่อนหนึ่งชั้น · เส้นแผนที่ปรับปรุงแล้วลดความชันช่วงนี้เหลือ 6.5 จุดต่อสัปดาห์ไว้แล้ว · ให้ย้ายงานที่ไม่ขึ้นกับผู้อื่นขึ้นไปทำใน Sprint 2 \\
กำลังคนฝั่ง Back-End ไม่พอเทียบสัดส่วนงาน & \textbf{สูง} & \textbf{สูง} & ในรอบนี้มีผู้คอมมิตโค้ดฝั่ง Back-End 2 คน ต่อน้ำหนักงาน 30\% ซึ่งสูงที่สุด · Procedure ที่ต้องเขียนใน Sprint 2 มีประมาณ 15 รายการจาก 60 รายการในสัญญา · ให้มอบหมาย Procedure ที่ไม่มีตรรกะซับซ้อนให้สมาชิกฝั่ง Front-End ช่วยเขียน \\
เชื่อมต่อครั้งแรกแล้วพบว่าโครงสร้างข้อมูลไม่ตรงกัน & ปานกลาง & \textbf{สูง} & ความเสี่ยงนี้ลดลงจากครั้งที่ 2 เพราะ \texttt{auth.me} พิสูจน์แล้วว่าเส้นทางจาก Prisma ถึง Zod Schema ทำงานได้ในไฟล์ทดสอบ · ยังคงเหลือความเสี่ยงจากรุ่น \texttt{zod} ที่ไม่ตรงกัน \\
หน้าจอที่ทำเสร็จแล้วต้องแก้ซ้ำเมื่อต่อข้อมูลจริง & \textbf{สูง} & ปานกลาง & จำนวนหน้าที่ต้องแก้เพิ่มจาก 7 เป็น 11 หน้าในรอบเดียว · ให้ต่อข้อมูลจริงกับหน้าโปรไฟล์เป็นหน้าแรกก่อนสร้างหน้าใหม่เพิ่ม \\
สคีมาถูกปรับใหญ่อีกครั้งหลังมีข้อมูลจริงแล้ว & ปานกลาง & \textbf{สูง} & สคีมาถูกปรับใหญ่ 2 ครั้งใน 8 วัน (12 และ 13 ส.ค.) ขณะที่ฐานข้อมูลยังว่าง · ให้ปิดรายการที่ยังค้าง 2 ข้อ (สถานะอนุมัติบางส่วน และคีย์ไม่ซ้ำ) ให้เสร็จก่อนเขียนสคริปต์ seed \\
\bottomrule
\end{tabularx}

% =============================================================================
\clearpage
\section{แผนการดำเนินงานที่ปรับปรุงแล้ว}
\label{sec:replan}

แผน Sprint ยึดตามที่เสนอไว้ในรายงานครั้งที่ 2 หัวข้อ 10 ไม่มีการเลื่อนเพิ่มในรอบนี้ หัวข้อย่อยแรกคือแผนสัปดาห์ถัดไปตามที่สไลด์ Week 06 หน้า 10 กำหนด

\subsection{แผนสัปดาห์ถัดไป (14 – 20 สิงหาคม 2569)}

\begin{tabularx}{\dimexpr\textwidth-\leftskip\relax}{@{} L{1.5cm} Y L{2.0cm} L{2.2cm} @{}}
\toprule
\hrow \thead{กำหนด} & \thead{งาน} & \thead{ผู้รับผิดชอบ} & \thead{ขึ้นกับ} \\
\midrule
14 ส.ค. & เพิ่มค่าอนุมัติบางส่วนใน \texttt{ApproveStatus} และ \texttt{@unique} บน \texttt{Email} กับ \texttt{UserID} แล้วสร้าง Migration ใหม่ & BE-lead + SA & --- \\
15 ส.ค. & ประชุมลงมติสัญญา แล้วคอมมิตไฟล์สัญญาพร้อมยกรุ่น \texttt{zod} ฝั่ง Front-End และถอด \texttt{continue-on-error} ใน CI ในคำขอรวมโค้ดเดียว & API + SA + FE-lead & งานวันที่ 14 ส.ค. \\
16 ส.ค. & สคริปต์ข้อมูลตั้งต้น 2 หน่วยงาน พร้อมผู้ใช้ครบทุกบทบาท & BE-dev & งานวันที่ 14 ส.ค. \\
17 ส.ค. & Procedure \texttt{auth.login} และ \texttt{auth.logout} พร้อมเติมการตรวจโทเคนใน \texttt{verifySessionToken()} & BE-dev & งานวันที่ 16 ส.ค. \\
18 ส.ค. & ติดตั้ง \texttt{vitest} ฝั่ง Front-End และย้าย \texttt{ac-checks.ts} เข้าเป็นชุดทดสอบ พร้อมเพิ่มงานตรวจชุดทดสอบเข้า CI & QA-lead + FE-lead & --- \\
19 ส.ค. & Procedure รายการอุปกรณ์ ค้นหา และจำนวนคงเหลือ & BE-dev & งานวันที่ 15 ส.ค. \\
20 ส.ค. & ต่อหน้าโปรไฟล์เข้ากับ \texttt{auth.me} ของจริง เป็นหน้าแรกที่เลิกใช้ข้อมูลจำลอง & FE-lead & งานวันที่ 17 ส.ค. \\
20 ส.ค. & \textbf{จุดตัดสินใจกลาง Sprint 2} ทบทวนว่าเส้นทางข้อมูลครบสามชั้นทำงานได้แล้วหรือไม่ & PM & --- \\
\bottomrule
\end{tabularx}

\subsection{Sprint 2 ``ยืม–คืนครบวงจร'' (14 – 27 สิงหาคม 2569)}

\begin{itemize}
  \item \textbf{ฝั่ง Back-End} ระบบเข้าสู่ระบบและควบคุมสิทธิ์ที่ใช้งานได้จริง · รายการอุปกรณ์และจำนวนคงเหลือ · สร้างคำขอ อนุมัติอัตโนมัติ ผูกชิ้น ส่งมอบ และรับคืน · สคริปต์ข้อมูลตั้งต้น · การจัดการข้อผิดพลาดและการตรวจความถูกต้อง
  \item \textbf{ฝั่ง Front-End} เปลี่ยนหน้าจอที่มีอยู่แล้ว 11 หน้าจากข้อมูลจำลองเป็นข้อมูลจริง · หน้าส่งคำขอและหน้าคิวของเจ้าหน้าที่ · ยกรุ่นไลบรารีตรวจสอบข้อมูล
  \item \textbf{ฝั่ง QA} แผนการทดสอบ ชุดทดสอบฝั่ง Front-End และการทดสอบเส้นทางปกติรวมกรณียืมข้ามภาควิชา
\end{itemize}

\noindent\textbf{เกณฑ์ ``เสร็จ''} ยืมอุปกรณ์ของภาควิชาคอมพิวเตอร์และภาควิชาไฟฟ้าในระบบเดียวได้ครบวงจรในเส้นทางปกติ และไม่เหลือข้อมูลจำลองในเส้นทางหลัก

\subsection{Sprint 3 ถึง Sprint 5}

ไม่มีการเปลี่ยนแปลงจากที่เสนอไว้ในรายงานครั้งที่ 2 คือ Sprint 3 (28 ส.ค. – 10 ก.ย.) รับเรื่องกฎธุรกิจและหลายหน่วยงานระดับที่ 1 · Sprint 4 (11 – 24 ก.ย.) รับเรื่องตะกร้าข้ามหน่วยงาน การจองแบบย่อ และงานขัดเงา โดยเป็น Sprint สุดท้ายที่เพิ่มความสามารถใหม่ได้ · Sprint 5 (25 ก.ย. – 1 ต.ค.) หยุดเพิ่มความสามารถ ทดสอบถดถอย นำระบบขึ้นใช้งานจริง และซ้อมนำเสนอ

\subsection{การประเมินกำหนดเสร็จของโครงการ}

คณะทำงานประเมินว่าโครงการยังเสร็จได้ภายในวันที่ 1 ตุลาคม 2569 ตามกำหนดเดิม โดยเงื่อนไขที่ต้องเป็นจริงมี 3 ข้อ

\begin{enumerate}
  \item ไฟล์สัญญาถูกคอมมิตเข้ารีโปพร้อมยกรุ่น \texttt{zod} ภายในวันที่ 15 สิงหาคม 2569
  \item เส้นทางข้อมูลครบสามชั้น คือหน้าจอเรียก tRPC แล้วอ่านจากฐานข้อมูลจริงได้อย่างน้อยหนึ่งหน้า ภายในวันที่ 20 สิงหาคม 2569
  \item วงจรยืม–คืนเส้นทางปกติทำงานได้จริงภายในวันที่ 27 สิงหาคม 2569 ซึ่งเป็นวันสิ้นสุด Sprint 2
\end{enumerate}

\noindent หากข้อใดข้อหนึ่งไม่เป็นไปตามนี้ คณะทำงานจะเสนอลดขอบเขตของ MVP ลงเหลือเฉพาะการยืมภายในหน่วยงานเดียว โดยตัดความสามารถข้ามหน่วยงานออก ซึ่งกระทบจุดขายหลักของโครงการโดยตรง จึงเสนอเป็นทางเลือกสุดท้าย

% =============================================================================
\section{ทรัพยากรและการสนับสนุนที่ต้องการเพิ่มเติม}

\begin{enumerate}
  \item \textbf{มติของที่ประชุมเรื่องสัญญา API} เป็นเงื่อนไขของงานฝั่ง Back-End ทุกรายการใน Sprint 2 งานเตรียมประกอบการตัดสินใจพร้อมแล้ว 64 หน้า 19 วาระ พร้อมแบบฟอร์มลงมติ
  \item \textbf{กำลังคนฝั่ง Back-End} ปัจจุบันมีผู้คอมมิตโค้ดฝั่ง Back-End 2 คน ต่อสัดส่วนน้ำหนักงาน 30\% ซึ่งสูงที่สุดในโครงการ และ Sprint 2 ต้องเขียน Procedure ประมาณ 15 รายการ ข้อเสนอนี้ยื่นไว้ตั้งแต่รายงานครั้งที่ 1
  \item \textbf{ฐานข้อมูล PostgreSQL สำหรับใช้ร่วมกันในทีม} ปัจจุบันแต่ละคนใช้ฐานข้อมูลบนเครื่องตนเอง ทำให้ทดสอบร่วมกันไม่ได้ และเป็นเหตุให้ไฟล์ทดสอบที่มีอยู่รันใน CI ไม่ได้ ค่าใช้จ่ายอยู่ในวงเงิน 1{,}500 บาทที่ประเมินไว้แล้ว
  \item \textbf{มติเรื่องการตัดขอบเขต 3 รายการ} ที่เสนอไว้ในรายงานครั้งที่ 2 และยังไม่มีมติ ซึ่งจำเป็นต่อการวางแผน Sprint 3 และ Sprint 4
  \item \textbf{คำแนะนำจากอาจารย์ที่ปรึกษาเรื่องการคงไว้ซึ่ง MongoDB} เป็นข้อค้างจากรายงานครั้งที่ 1 ที่ยังไม่มีข้อสรุป
\end{enumerate}

% =============================================================================
\section{ความสำเร็จที่บันทึกไว้ในรอบนี้}

\begin{itemize}
  \item \textbf{tRPC Procedure รายการแรกเข้าสาย \texttt{main} แล้ว} มิติ Back-End Dev ซึ่งอยู่ที่ 12\% เท่าเดิมมาสองรอบรายงาน ขยับเป็น 20\% ในรอบนี้ พร้อมชั้นควบคุมสิทธิ์ที่ครอบคลุมทั้ง 4 บทบาทที่ระบบต้องใช้
  \item \textbf{ไฟล์ทดสอบชุดแรกของโครงการ} มิติ Test ขยับจาก 0\% เป็นค่าบวกเป็นครั้งแรกหลังผ่านไปสองรอบรายงาน
  \item \textbf{ระบบตรวจสอบอัตโนมัติทำงานบนสายหลักแล้ว} run 2 ครั้งล่าสุดผ่านทั้งคู่ ใช้เวลา 34 และ 40 วินาที คำขอรวมโค้ดทุกรายการนับจากนี้จะถูกตรวจการติดตั้ง ชนิดข้อมูล และการ Build ของทั้งสองฝั่งโดยอัตโนมัติ
  \item \textbf{หน้าจอฝั่งผู้ยืมชุดแรกเสร็จก่อนกำหนด} หน้ารายการอุปกรณ์ รายละเอียดอุปกรณ์ และรายการห้อง รวม 1{,}512 บรรทัด เป็นงานที่แผนปรับปรุงกำหนดไว้ใน Sprint 2
  \item \textbf{ข้อค้างเรื่องตารางผูกชิ้นอุปกรณ์ถูกปิด} เป็นรายการเดียวจากสามรายการที่ค้างมาตั้งแต่รายงานครั้งที่ 1 ที่ได้รับการแก้ไขแล้ว
\end{itemize}

% =============================================================================
\clearpage
\section*{ภาคผนวก ก — ที่มาของตัวเลขความก้าวหน้า}
\addcontentsline{toc}{section}{ภาคผนวก ก — ที่มาของตัวเลขความก้าวหน้า}

ตัวเลขในคอลัมน์ ``จริง'' ของหัวข้อ \ref{sec:dimension} ประเมินจากหลักฐานในรีโปที่ commit \headsha\ ณ วันที่ \reportdate\ โดยไม่นับโฟลเดอร์ \texttt{spike-option-b/} และไม่นับโค้ดที่ Prisma สร้างอัตโนมัติ

\begin{tabularx}{\dimexpr\textwidth-\leftskip\relax}{@{} L{2.6cm} L{1.0cm} Y @{}}
\toprule
\hrow \thead{มิติงาน} & \thead{ผลจริง} & \thead{หลักฐานและวิธีคำนวณ} \\
\midrule
Back-End Design & 64\% &
ให้น้ำหนักแบบจำลองข้อมูลและสัญญา API อย่างละครึ่ง \newline
\textbf{แบบจำลองข้อมูล 88\%} จาก \texttt{schema.prisma} 27 ตาราง 9 Enum และ Migration \texttt{20260813063940\zb\_init} ที่สร้างตารางได้จริง เพิ่มจาก 85\% ในรอบที่แล้วเพราะข้อค้างเรื่องตารางผูกชิ้นอุปกรณ์ถูกปิด หักส่วนที่ยังขาดคือค่าสถานะอนุมัติบางส่วนและคีย์ไม่ซ้ำบนอีเมล \newline
\textbf{สัญญา API 40\%} เพิ่มจาก 25\% ในรอบที่แล้ว เพราะมีสัญญาที่เป็นโค้ดจริงแล้วทั้งสองฝั่ง คือ 60 Procedure ฝั่ง Front-End และ Zod Schema กับ 1 Procedure ฝั่ง Back-End หักส่วนที่ยังขาดคือไม่มีแหล่งความจริงเดียว ยังไม่มีมติที่ประชุม และรุ่น \texttt{zod} ไม่ตรงกัน \newline
$(88+40)/2 = 64$ \\

Front-End Design & 76\% &
ระบบออกแบบและเปลือกระบบ (น้ำหนัก 40\%) เสร็จ 100\% · แบบหน้าจอ (น้ำหนัก 35\%) มีเนื้อหาจริง 11 จาก 28 หน้า คิดเป็น 39\% · การออกแบบกระบวนการทางธุรกิจ (น้ำหนัก 25\%) เสร็จ 90\% \newline
$0.40(100)+0.35(39)+0.25(90) = 76$ \\

Back-End Dev & 20\% &
แจกแจงเป็นสี่ส่วนตามลำดับที่ต้องทำ \newline
\textbf{โครงสร้างพื้นฐานและ bootstrap (น้ำหนัก 12\%)} เสร็จ 100\% — ค่านี้เท่ากับ 12\% ที่รายงานไว้ในครั้งที่ 1 และ 2 พอดี จึงเทียบสามรอบได้ตรง \newline
\textbf{ชั้นยืนยันตัวตนและควบคุมสิทธิ์ (น้ำหนัก 18\%)} เสร็จ 40\% — มิดเดิลแวร์ 4 บทบาทและโครง Context ครบ แต่ \texttt{verifySessionToken()} ยังเป็นค่าคงที่ \texttt{null} และไม่มี Procedure เข้าสู่ระบบ \newline
\textbf{Procedure ตามสัญญา (น้ำหนัก 55\%)} เสร็จ 1.7\% — 1 จาก 60 รายการ \newline
\textbf{ข้อมูลตั้งต้น (น้ำหนัก 15\%)} เสร็จ 0\% \newline
$0.12(100)+0.18(40)+0.55(1.7)+0.15(0) = 20.1$ \\

Front-End Dev & 46\% &
ให้น้ำหนักหน้าจอ 70\% และโครงสร้างพื้นฐาน 30\% \newline
\textbf{หน้าจอ 27.5\%} จาก 28 หน้า มีเนื้อหาจริง 11 หน้า ทุกหน้ายังอ่านข้อมูลจำลอง จึงคิดให้หน้าละ 0.7 ของหน้าที่ต่อข้อมูลจริงแล้ว ได้ $11\times0.7 = 7.7$ เทียบ 28 หน้า คิดเป็น 27.5\% (ใช้ค่าปรับลด 0.7 ค่าเดียวกับรายงานครั้งที่ 2) \newline
\textbf{โครงสร้างพื้นฐาน 88\%} เพิ่มจาก 82\% เพราะ CI เข้าสายหลักแล้วและมีชั้นเชื่อมต่อ tRPC พร้อม \texttt{TRPCProvider} หักส่วนที่ยังไม่มีผลคือยังไม่มีไฟล์ใดเรียก \texttt{useTRPC()} \newline
$0.70(27.5)+0.30(88) = 45.7$ \\

Test & 10\% &
แจกแจงเป็นสามส่วน \newline
\textbf{ชุดทดสอบฝั่ง Back-End (น้ำหนัก 40\%)} เสร็จ 8\% — \texttt{auth.service.spec.ts} 99 บรรทัด 2 กรณี ครอบคลุม 1 จาก 60 Procedure และต้องต่อฐานข้อมูลจริงจึงยังรันใน CI ไม่ได้ ไม่นับ \texttt{app.controller.spec.ts} กับ \texttt{prisma.service.spec.ts} ซึ่งเป็นไฟล์ตัวอย่างที่มากับโครงเริ่มต้นของ NestJS \newline
\textbf{ชุดทดสอบฝั่ง Front-End (น้ำหนัก 40\%)} เสร็จ 10\% — ยังไม่มีตัวรันชุดทดสอบติดตั้ง มีเพียง \texttt{ac-checks.ts} 116 บรรทัดที่รันด้วยมือ \newline
\textbf{แผนทดสอบและเกณฑ์การยอมรับ (น้ำหนัก 20\%)} เสร็จ 15\% — มีเกณฑ์ยอมรับข้อ 2.4 ถึง 2.7 เป็นโค้ดตรวจได้ แต่ยังไม่มีเอกสารแผนการทดสอบ \newline
$0.40(8)+0.40(10)+0.20(15) = 10.2$ \\

Documentation & 64\% &
เพิ่มจาก 62\% ในรอบที่แล้ว จากแบบฟอร์มลงมติสัญญาที่เพิ่มเข้ามาในรอบนี้ รวมกับข้อเสนอโครงการ 42 หน้า คู่มือสัญญา 31 หน้า วาระประชุม 33 หน้า รายงานผลงานทดลอง 31 หน้า สารบัญไฟล์ และรายงานความก้าวหน้า 2 ฉบับที่มีอยู่เดิม \newline
ส่วนที่ยังขาดคือ SRS ฉบับทางการ SDS นอกเหนือบทฐานข้อมูล คู่มือผู้ใช้ และคู่มือติดตั้ง \\
\bottomrule
\end{tabularx}

\vspace{1em}
\noindent\textbf{การรวมเป็นค่าเดียว}
$0.10(64)+0.10(76)+0.30(20)+0.25(46)+0.15(10)+0.10(64)
= 6.4+7.6+6.0+11.5+1.5+6.4 = 39.4$ ปัดเป็น \actualpct

\vspace{0.8em}
\noindent\textbf{ข้อจำกัดของการประเมิน} ตัวเลขข้างต้นเป็นการแปลงปริมาณผลงานที่ปรากฏในรีโปให้เป็นร้อยละ ซึ่งย่อมมีดุลพินิจของผู้ประเมินอยู่ในสามจุด คือค่าน้ำหนักของแต่ละมิติ · ค่าปรับลด 0.7 ที่ใช้กับหน้าจอที่ยังอ่านข้อมูลจำลอง · และการแจกแจงน้ำหนักย่อยภายในมิติ Back-End Dev กับมิติ Test ซึ่งรอบนี้ระบุออกมาเป็นตัวเลขเป็นครั้งแรก คณะทำงานจึงรายงานทั้งค่าน้ำหนักและวิธีคำนวณไว้ครบ เพื่อให้ผู้อ่านปรับค่าตามที่เห็นสมควรได้ ค่าน้ำหนักย่อยของมิติ Back-End Dev ถูกเลือกให้ส่วน ``โครงสร้างพื้นฐาน'' เท่ากับ 12\% พอดี เพื่อให้ค่าที่รายงานไว้ในครั้งที่ 1 และ 2 ยังเทียบกับรอบนี้ได้โดยตรง

% =============================================================================
\clearpage
\section*{ภาคผนวก ข — ค่าที่ใช้พล็อตกราฟ S-Curve}
\addcontentsline{toc}{section}{ภาคผนวก ข — ค่าที่ใช้พล็อตกราฟ S-Curve}

\subsection*{ข.1 เส้นแผนที่ปรับปรุงแล้ว}

ค่าเป้าหมายรายมิติ ณ วันสิ้นสุดของแต่ละ Sprint คูณด้วยน้ำหนักของมิตินั้นแล้วรวมกัน คอลัมน์ ``รวม'' คือค่าที่พล็อตบนกราฟ

\begin{tabularx}{\dimexpr\textwidth-\leftskip\relax}{@{} L{2.4cm} L{1.0cm} L{1.0cm} L{1.0cm} L{1.0cm} L{1.0cm} L{1.0cm} L{1.0cm} Y @{}}
\toprule
\hrow \thead{ขอบ Sprint} & \thead{สัปดาห์} & \thead{BE Des} & \thead{FE Des} & \thead{BE Dev} & \thead{FE Dev} & \thead{Test} & \thead{Doc} & \thead{รวม} \\
\hrow \thead{} & \thead{} & \thead{10\%} & \thead{10\%} & \thead{30\%} & \thead{25\%} & \thead{15\%} & \thead{10\%} & \thead{100\%} \\
\midrule
เริ่ม 24 ก.ค. & 0.00 & 0 & 0 & 0 & 0 & 0 & 10 & \textbf{1} \\
S0 · 30 ก.ค. & 0.86 & 45 & 30 & 5 & 5 & 0 & 20 & \textbf{12} \\
S1 · 13 ส.ค. & 2.86 & 88 & 75 & 25 & 42 & 10 & 65 & \textbf{42} \\
S2 · 27 ส.ค. & 4.86 & 95 & 85 & 60 & 65 & 35 & 70 & \textbf{65} \\
S3 · 10 ก.ย. & 6.86 & 100 & 92 & 78 & 78 & 52 & 76 & \textbf{78} \\
S4 · 24 ก.ย. & 8.86 & 100 & 100 & 97 & 97 & 80 & 88 & \textbf{94} \\
S5 · 1 ต.ค. & 9.86 & 100 & 100 & 100 & 100 & 100 & 100 & \textbf{100} \\
\bottomrule
\end{tabularx}

\vspace{0.8em}
\noindent ตัวอย่างการคำนวณแถว S1 · 13 ส.ค. ซึ่งเป็นวันที่ออกรายงานฉบับนี้
$0.10(88)+0.10(75)+0.30(25)+0.25(42)+0.15(10)+0.10(65)
= 8.8+7.5+7.5+10.5+1.5+6.5 = 42.3$ ปัดเป็น 42

\vspace{0.8em}
\noindent ความชันของเส้นแผนระหว่างขอบ Sprint คำนวณจากผลต่างของคอลัมน์ ``รวม'' หารด้วยจำนวนสัปดาห์ของช่วงนั้น

\begin{tabularx}{\dimexpr\textwidth-\leftskip\relax}{@{} L{1.6cm} L{2.3cm} L{2.2cm} L{3.1cm} Y @{}}
\toprule
\hrow \thead{ช่วง} & \thead{ระยะเวลา} & \thead{ความก้าวหน้า} & \thead{ความชัน} & \thead{หมายเหตุ} \\
\midrule
Sprint 0 & 0.86 สัปดาห์ & $1 \rightarrow 12$ & 12.8 จุด/สัปดาห์ & งานออกแบบและติดตั้งโครงสร้าง \\
Sprint 1 & 2.00 สัปดาห์ & $12 \rightarrow 42$ & \textbf{15.0 จุด/สัปดาห์} & ชันที่สุด \\
Sprint 2 & 2.00 สัปดาห์ & $42 \rightarrow 65$ & 11.5 จุด/สัปดาห์ & \\
Sprint 3 & 2.00 สัปดาห์ & $65 \rightarrow 78$ & \textbf{6.5 จุด/สัปดาห์} & ทับช่วงสอบกลางภาค \\
Sprint 4 & 2.00 สัปดาห์ & $78 \rightarrow 94$ & 8.0 จุด/สัปดาห์ & \\
Sprint 5 & 1.00 สัปดาห์ & $94 \rightarrow 100$ & 6.0 จุด/สัปดาห์ & หยุดเพิ่มความสามารถ \\
\bottomrule
\end{tabularx}

\subsection*{ข.2 เส้นแผนเดิม}

ชุดพิกัดที่ใช้ในรายงานครั้งที่ 1 และ 2 พล็อตไว้ในกราฟเพื่อการเปรียบเทียบ

\begin{tabularx}{\dimexpr\textwidth-\leftskip\relax}{@{} L{2.0cm} L{1.6cm} L{1.6cm} L{1.6cm} L{1.6cm} L{1.6cm} L{1.6cm} Y @{}}
\toprule
\hrow \thead{สัปดาห์} & \thead{0} & \thead{1} & \thead{3} & \thead{5} & \thead{7} & \thead{9} & \thead{10} \\
\midrule
\% สะสม & 0 & 25 & 55 & 75 & 90 & 97 & 100 \\
\bottomrule
\end{tabularx}

\vspace{0.8em}
\noindent ค่า ณ สัปดาห์ที่ 2.86 ซึ่งเป็นวันที่ออกรายงานฉบับนี้ ได้จากการประมาณเชิงเส้นระหว่างสัปดาห์ที่ 1 กับสัปดาห์ที่ 3 คือ $25 + (1.86/2)(55-25) = 52.9$ ปัดเป็น \planbasepct

\subsection*{ข.3 เส้นผลจริง}

\begin{tabularx}{\dimexpr\textwidth-\leftskip\relax}{@{} L{2.6cm} L{1.8cm} L{1.4cm} L{2.6cm} Y @{}}
\toprule
\hrow \thead{จุด} & \thead{วันที่} & \thead{สัปดาห์} & \thead{\% สะสม} & \thead{ที่มา} \\
\midrule
เริ่มโครงการ & 24 ก.ค. 2569 & 0.00 & 0 & --- \\
รายงานครั้งที่ 1 & 7 ส.ค. 2569 & 2.00 & 24 & \texttt{docs/build/progress.pdf} \\
รายงานครั้งที่ 2 & 10 ส.ค. 2569 & 2.43 & 32 & \texttt{docs/build/progress2.pdf} \\
รายงานครั้งที่ 3 & 13 ส.ค. 2569 & 2.86 & \textbf{39} & ฉบับนี้ · ภาคผนวก ก \\
\bottomrule
\end{tabularx}

\vspace{0.8em}
\noindent ทั้งสามจุดใช้ค่าน้ำหนักมิติงานชุดเดียวกันและวิธีคำนวณเดียวกัน จึงเทียบกันได้โดยตรง

% =============================================================================
\clearpage
\section*{ภาคผนวก ค — หลักฐานอ้างอิงในรีโป}
\addcontentsline{toc}{section}{ภาคผนวก ค — หลักฐานอ้างอิงในรีโป}

รายการต่อไปนี้คือหลักฐานที่ใช้ยืนยันข้อความสถานะทุกข้อในรายงานฉบับนี้ ตรวจสอบได้ที่ \texttt{origin/main} commit \headsha

\begin{tabularx}{\dimexpr\textwidth-\leftskip\relax}{@{} L{5.4cm} Y @{}}
\toprule
\hrow \thead{ข้อความในรายงาน} & \thead{หลักฐาน} \\
\midrule
มี Procedure 1 รายการ & \texttt{backend/src/auth/auth.router.ts} — \texttt{@Router(\{ alias: 'auth' \})} มี \texttt{@Query} หนึ่งรายการคือ \texttt{me()} \\
Procedure เรียกสำเร็จไม่ได้ & \texttt{backend/src/trpc/context.ts} — \texttt{verifySessionToken()} คืน \texttt{null} เสมอ พร้อมคอมเมนต์ \texttt{TODO} \\
ชั้นควบคุมสิทธิ์ 4 บทบาท & \texttt{backend/src/trpc/auth.middleware.ts} — \texttt{AuthMiddleware} และ \texttt{requireRole()} ที่สร้าง \texttt{StaffMiddleware}, \texttt{Supervisor\zb Middleware}, \texttt{AdminMiddleware} · ทั้งหมดลงทะเบียนใน \texttt{app.module.ts} \\
ไฟล์ทดสอบชุดแรก & \texttt{backend/src/auth/auth.service.spec.ts} — 99 บรรทัด 2 บล็อก \texttt{it()} \\
สคีมา 27 ตาราง 9 Enum & \texttt{backend/prisma/schema.prisma} — นับด้วย \texttt{grep -c '\^{}model '} ได้ 27 และนับ \texttt{\^{}enum } ได้ 9 \\
Migration ล่าสุด & \texttt{backend/prisma/migrations/20260813063940\zb\_init/migration.sql} — 538 บรรทัด \\
ไม่มีสคริปต์ข้อมูลตั้งต้น & ไม่มีไฟล์ seed ใน \texttt{backend/prisma/} และไม่มีคีย์ \texttt{prisma.seed} ใน \texttt{backend/package.json} \\
\texttt{Email} ไม่มี \texttt{@unique} & \texttt{schema.prisma} model \texttt{AccountInfo} — ฟิลด์ \texttt{Email} และ \texttt{UserID} ประกาศเป็น \texttt{String} เปล่า \\
\texttt{ApproveStatus} มี 4 ค่า & \texttt{schema.prisma} — \texttt{Pending}, \texttt{Approved}, \texttt{Rejected}, \texttt{Canceled} \\
หน้าจอมีเนื้อหาจริง 11 จาก 28 & นับไฟล์ \texttt{frontend/src/features/**/*-page.tsx} ได้ 28 ไฟล์ · 17 ไฟล์มี 6 บรรทัด (หน้าเปล่า) · 11 ไฟล์เหลือมี 124 ถึง 612 บรรทัด \\
เส้นทางหน้าจอ 28 เส้นทาง & \texttt{frontend/src/app/router.tsx} — \texttt{<Route path=} 29 รายการ โดย 1 รายการเป็นเส้นทางสำรอง \texttt{path="*"} \\
ไม่มีไฟล์ใดเรียก \texttt{useTRPC()} & ค้นทั้ง \texttt{frontend/src/} พบ \texttt{useTRPC} เฉพาะในไฟล์ที่นิยามมันคือ \texttt{lib/trpc.ts} \\
สัญญาฝั่ง Front-End 60 Procedure & \texttt{frontend/src/server/trpc-contract.ts} — 188 บรรทัด นับ \texttt{proc.query} กับ \texttt{proc.\dots mutation} รวมได้ 60 รายการใน 11 กลุ่ม \\
รุ่น \texttt{zod} ไม่ตรงกัน & \texttt{backend/package.json} ระบุ \texttt{"zod": "\^{}4.4.3"} · \texttt{frontend/package.json} ระบุ \texttt{"zod": "\^{}3.23.8"} \\
สคริปต์ตรวจเวอร์ชันรายงานปัญหา 3 ข้อ & รัน \texttt{node scripts/check-versions.mjs} ได้ผล 3 FAIL และ exit code 1 \\
CI อยู่บนสายหลักและผ่าน & \texttt{.github/workflows/ci.yml} อยู่บน \texttt{main} ตั้งแต่ PR \#19 · run ล่าสุด 2 รายการบน \texttt{main} และบนคำขอรวมโค้ด สถานะ success เมื่อ 13 ส.ค. \\
งานตรวจเวอร์ชันไม่บังคับใช้ & \texttt{ci.yml} งาน \texttt{versions} มี \texttt{continue-on-error: true} \\
ปริมาณงานในรอบ \deltaperiod & \texttt{git diff --shortstat 53a42eb e76f01f} เฉพาะพาธที่ไม่ใช่โค้ดที่สร้างอัตโนมัติ ได้ 53 ไฟล์ เพิ่ม 5{,}365 บรรทัด ลบ 292 บรรทัด \\
ผู้คอมมิตฝั่ง Back-End 2 คน & \texttt{git log --since=2026-08-10 -- backend/} มีผู้เขียนคอมมิต 2 ราย \\
\bottomrule
\end{tabularx}

\end{document}
